\chapter{Pseudodifferential operators on manifolds}\label{ch:manifolds}
\providecommand{\Gsm}{G^{-\infty}}
\providecommand{\Symq}[1]{S^{#1-[1]}}

In this chapter pseudodifferential operators are transferred from $\R^n$ to a
compact manifold $X$, acting between sections of vector bundles. Two definitions
are given and shown to agree: the invariant one, in which an operator of order
$m$ is identified with its Schwartz kernel, a distribution on $X^2$ conormal of
order $m$ to the diagonal (using the conormal distributions of
\cref{ch:conormal}); and the local one, in which an operator is required to be,
in every coordinate patch, an element of the calculus on $\R^n$ of
\cref{ch:euclidean}, and to have a kernel smooth away from the diagonal. With
the coordinate invariance established in \cref{ch:euclidean,ch:microlocal},
the whole symbol calculus carries over. We then construct parametrices for
elliptic operators, introduce Sobolev spaces on $X$, and prove that elliptic
operators are Fredholm, with generalized inverses which are again
pseudodifferential. Smoothing operators, which carry all the ``errors'' of the
calculus, are discussed in some detail: they form an ideal of ``infinite
matrices'' with a trace and a group of invertible elements, $\Gsm$, which is
the starting point for the index theory of later chapters. The main
applications here are the Hodge theorem for elliptic complexes, in particular
for the de~Rham complex, and the existence of pseudodifferential projections
onto prescribed subspaces. The facts about bounded operators on Hilbert space
which are used are collected, with proofs, in \cref{ch:hilbert}.

Throughout, unless explicitly stated otherwise, $X$ is a \emph{compact}
manifold of dimension $n$; several definitions extend to the non-compact case
with minor changes, which are indicated in remarks.

\section{Smooth manifolds}\label{manifolds:sec:manifolds}

\subsection{$\CI$ structures}
Let $X$ be a paracompact Hausdorff topological space. We describe the smooth
structure of $X$ through its algebra of smooth functions.

\begin{definition}\label{manifolds:def:structure}
A \emph{$\CI$ structure} on $X$ is a subspace $\mathcal F\subset\Cz(X;\R)$
with the following property: for each $\bar x\in X$ there exist
$f_1,\dots,f_n\in\mathcal F$ and an open neighbourhood $\Omega\ni\bar x$ such
that
\begin{equation}\label{manifolds:eq:coord}
\kappa:\Omega\ni x\longmapsto (f_1(x),\dots,f_n(x))\in\R^n
\end{equation}
is a homeomorphism onto an open subset of $\R^n$ and every $f\in\mathcal F$
satisfies $f|_{\Omega}=g\circ\kappa$ for some $g\in\CI(\kappa(\Omega))$. The map
$\kappa$ is then called a \emph{coordinate system} (or chart) near $\bar x$.
Two $\CI$ structures $\mathcal F_1$, $\mathcal F_2$ are \emph{compatible} if
$\mathcal F_1\cup\mathcal F_2$ is also a $\CI$ structure.
\end{definition}

Compatibility is an equivalence relation, and each equivalence class contains a
unique maximal element, the union of the class (\cref{manifolds:ex:compatible}).

\begin{definition}\label{manifolds:def:manifold}
A \emph{$\CI$ manifold} is a paracompact Hausdorff space $X$ with a maximal
$\CI$ structure, which is denoted $\CI(X)\subset\Cz(X)$.
\end{definition}

The maximal structure is an algebra (\cref{manifolds:ex:maximal}), and the
number $n$ of coordinate functions is locally constant; we always assume it
is constant, $n=\dim X$. If $\kappa:\Omega\to\kappa(\Omega)$ is a coordinate
system in $X$ and $\CIc(\Omega)\subset\CI(X)$ denotes the space of smooth
functions vanishing outside a compact subset of $\Omega$, then
\begin{equation*}
\kappa^*:\CIc(\kappa(\Omega))\longrightarrow\CIc(\Omega)
=\{u\in\CI(X);\ u=0 \text{ on } X\setminus K,\ K\Subset\Omega\}
\end{equation*}
is an isomorphism. Moreover $\CI(X)$ is \emph{local}: if $u:X\to\R$ and each
$\bar x\in X$ has an open neighbourhood $\Omega_{\bar x}$ on which
$u=f_{\bar x}$ for some $f_{\bar x}\in\CI(X)$, then $u\in\CI(X)$. The transition
maps $\kappa'\circ\kappa^{-1}$ between two coordinate systems are
diffeomorphisms between open subsets of $\R^n$, so this definition agrees with
the usual one by atlases. A map $G:X\to Y$ between manifolds is $\CI$ if
$G^*\CI(Y)\subset\CI(X)$, that is, $u\circ G\in\CI(X)$ for all $u\in\CI(Y)$.
Complex-valued smooth functions are those with real and imaginary parts in
$\CI(X)$.

\subsection{Partitions of unity}
Any $\CI$ manifold admits partitions of unity subordinate to any open cover
(\cref{manifolds:ex:partition}). For compact manifolds this is elementary: for
each point $p$ choose a coordinate system near $p$ inside an element of the
cover and pull back a non-negative bump function on $\R^n$ which is positive at
the image of $p$; by locality this gives an element of $\CI(X)$. Finitely many
of these are positive on a set covering $X$, and dividing each of them by
their (positive) sum gives a finite partition of unity subordinate to the cover.
For localizing operators a slightly stronger property is convenient.

\begin{lemma}\label{manifolds:lem:good-partition}
If $\{U_i\}$ is an open cover of a compact manifold $X$ there is a finite
partition of unity $\{\phi_a\}$, $0\le\phi_a\le1$, $\sum_a\phi_a=1$, each
$\supp\phi_a$ contained in some $U_i$, such that
\begin{equation}\label{manifolds:eq:good-partition}
\supp\phi_a\cap\supp\phi_b\ne\emptyset\Longrightarrow
\supp\phi_a\cup\supp\phi_b\subset U_i\ \text{ for some } i.
\end{equation}
\end{lemma}
\begin{proof}
Choose a finite partition of unity $\{\chi_i\}$ with $\supp\chi_i\subset U_i$
(passing to a finite subcover). For $p\in X$ set
\begin{equation*}
V_p=\bigcap_{i:\,p\in\supp\chi_i}U_i\ \setminus\bigcup_{i:\,p\notin\supp\chi_i}\supp\chi_i ,
\end{equation*}
an open neighbourhood of $p$ with the property that $V_p\cap\supp\chi_i\ne\emptyset$
implies $p\in\supp\chi_i$ and hence $V_p\subset U_i$. Choose finitely many
$p_j$ such that the $V_j=V_{p_j}$ cover $X$, a partition of unity $\{\psi_j\}$
with $\supp\psi_j\subset V_j$, and set $\phi_{ij}=\chi_i\psi_j$. If
$\supp\phi_{ij}\cap\supp\phi_{i'j'}\ne\emptyset$ then $\supp\psi_j$ meets
$\supp\chi_{i'}$, so $V_j\subset U_{i'}$, and therefore
$\supp\phi_{ij}\subset U_{i'}$; since also $\supp\phi_{i'j'}\subset\supp\chi_{i'}\subset U_{i'}$,
the union of the two supports lies in $U_{i'}$.
\end{proof}

\section{Vector bundles, densities and distributions}\label{manifolds:sec:bundles}

\subsection{Vector bundles}
A (complex) \emph{vector bundle} of rank $N$ over $X$ is a manifold $E$ with a
smooth surjection $\pi:E\to X$, each fibre $E_x=\pi^{-1}(x)$ carrying the
structure of an $N$-dimensional complex vector space, such that each point of
$X$ has an open neighbourhood $U$ and a diffeomorphism
$\pi^{-1}(U)\to U\times\C^N$ which is fibrewise linear and covers the identity
(a \emph{local trivialization}). Two trivializations differ by a smooth map
$U\cap U'\to\GL(N,\C)$. Real vector bundles are defined in the same way. The
space of smooth sections is
\begin{equation*}
\CI(X;E)=\{u:X\to E \text{ smooth};\ \pi\circ u=\Id\};
\end{equation*}
it is a module over $\CI(X)$, and over a trivializing open set a section is just
an $N$-tuple of smooth functions. We shall often use charts $\kappa:W\to\kappa(W)$
over which the bundles under consideration are trivial; we call these
\emph{trivializing charts}. Since $X$ is compact, it is covered by finitely many
of them.

\subsubsection*{The cotangent bundle}
For $\bar x\in X$ let $\mathcal I_{\bar x}=\{u\in\CI(X);u(\bar x)=0\}$ and
let $\mathcal I_{\bar x}^2$ be the space of finite sums of products of pairs
of elements of $\mathcal I_{\bar x}$. The cotangent space is
\begin{equation*}
T^*_{\bar x}X=\mathcal I_{\bar x}/\mathcal I^2_{\bar x},\qquad T^*X=\bigcup_{\bar x\in X}T^*_{\bar x}X,
\end{equation*}
and the differential of $f\in\CI(X)$ is the map $df:X\to T^*X$,
$df(\bar x)=[f-f(\bar x)]\in T^*_{\bar x}X$. In a coordinate system $x_1,\dots,x_n$
near $\bar x$, Taylor's formula with integral remainder shows that
$f-f(\bar x)-\sum_j\pa_jf(\bar x)(x_j-\bar x_j)\in\mathcal I^2_{\bar x}$ (after
multiplication by a cutoff equal to one near $\bar x$), so $dx_1,\dots,dx_n$ is a
basis of $T^*_{\bar x}X$ and $df=\sum_j\pa_jf\,dx_j$. The coordinates
$(x,\xi)\mapsto\sum_j\xi_j\,dx_j$ make $T^*X$ a real vector bundle; if $y=y(x)$
is another coordinate system, the corresponding fibre coordinates are related by
\begin{equation}\label{manifolds:eq:cotangent-change}
\xi=\Big(\frac{\pa y}{\pa x}\Big)^{T}\eta,\quad\text{i.e. }
\sum_j\xi_j\,dx_j=\sum_k\eta_k\,dy_k ,
\end{equation}
which is linear in the fibre variables with smooth invertible coefficients.
Each $df$ is a smooth section of $T^*X$.

\subsubsection*{Operations on bundles}
Any functorial operation on finite-dimensional vector spaces produces new
bundles from old, the transition functions being transformed accordingly. In
particular we have the tangent bundle $T_xX=(T^*_xX)^*$, the dual bundle $E'$,
the bundles $E\otimes F$, $\Hom(E,F)=F\otimes E'$, $\End(E)=\Hom(E,E)$, and the
exterior powers
\begin{equation*}
\Lambda^k_xX=\{u:(T_xX)^k\to\C\ \text{multilinear and antisymmetric}\},
\qquad \Lambda^1X=T^*X\otimes\C,
\end{equation*}
whose sections are the $k$-forms. If $E\to Y$ is a vector bundle and
$G:X\to Y$ is smooth, the \emph{pull-back bundle} $G^*E=\bigcup_{x\in X}E_{G(x)}$
is a smooth vector bundle over $X$; if $\phi\in\CI(Y;E)$ then $G^*\phi\in\CI(X;G^*E)$
and $\CI(X;G^*E)$ is spanned over $\CI(X)$ by $G^*\CI(Y;E)$. If $E\to X$ and
$F\to Y$ are bundles, the \emph{exterior tensor product} $F\boxtimes E\to X\times Y$
has fibre $(F\boxtimes E)_{(x,y)}=F_x\otimes E_y$. On $X^2$ we use the projections
$\pi_L(x,y)=x$ and $\pi_R(x,y)=y$, and the ``big'' homomorphism bundle
\begin{equation}\label{manifolds:eq:big-hom}
\Hom(E,F)=\pi_L^*F\otimes\pi_R^*E'=F\boxtimes E',\qquad \Hom(E,F)_{(x,y)}=\hom(E_y,F_x),
\end{equation}
whose restriction to the diagonal is $\hom(E,F)$. A \emph{hermitian metric} on
$E$ is a smoothly varying inner product on the fibres; one exists by a partition
of unity argument.

\subsection{Densities and integration}
For $\alpha\in\R$ the bundle of $\alpha$-densities is
\begin{equation*}
\Omega^\alpha_xX=\{u:\Lambda^n T_xX\setminus 0\to\C;\ u(t\,v)=|t|^\alpha u(v)\ \forall t\in\R\setminus0\},
\end{equation*}
a line bundle, and $\Omega=\Omega^1$. In local coordinates an $\alpha$-density
is $u=g(x)|dx|^\alpha$, and under a change of coordinates $y=y(x)$ the
coefficient transforms by $g\mapsto g\,|\det(\pa y/\pa x)|^{\alpha}$. Thus, by the
change of variables formula, if $u\in\CIc(W;\Omega)$ is supported in a coordinate
patch then $\int u=\int_{\R^n}g(x)\,dx$ is independent of the coordinates, and
using a partition of unity the integral
\begin{equation*}
\int_X:\CI(X;\Omega)\longrightarrow\C
\end{equation*}
is well defined. Densities, unlike forms, require no orientation. A \emph{smooth
positive density} $\nu$ (one which is a positive multiple of $|dx|$ in every
coordinate system) exists, again by a partition of unity, and any other density is
$f\nu$ with $f\in\CI(X)$. Given $\nu$ and a hermitian metric on $E$, the space
$L^2(X;E)$ is the completion of $\CI(X;E)$ with respect to
\begin{equation}\label{manifolds:eq:L2}
\inner uv=\int_X\langle u(x),v(x)\rangle_{E_x}\,\nu(x).
\end{equation}
Different choices give the same space with equivalent norms. On $X^2$ we use the
\emph{right density bundle} $\Omega_R=\pi_R^*\Omega$; a section of
$\Hom(E,F)\otimes\Omega_R$ is in local coordinates and trivializations a matrix
$k(x,y)|dy|$.

\subsection{Distributional sections}
The topology on $\CI(X;E)$ is given by the seminorms
$\sup|D^\alpha(\kappa_i^{-1})^*(\phi_iu)|$, $\alpha\in\N_0^n$, for a finite cover
by trivializing charts $\kappa_i$ and a subordinate partition of unity $\phi_i$
(the $N$-tuples of functions being estimated componentwise); this makes
$\CI(X;E)$ a Fr\'echet space and is independent of choices. The space of
distributional sections is the dual space
\begin{equation}\label{manifolds:eq:distributions}
\CmI(X;E)=\big(\CI(X;E'\otimes\Omega)\big)',
\end{equation}
with the weak topology. Integration of the pairing gives the inclusion
$\CI(X;E)\hookrightarrow\CmI(X;E)$, $u\mapsto(\phi\mapsto\int_X\phi(u))$, and
$\CI(X;E)$ is dense. Multiplication by $\CI(X)$ and the restriction to open sets
are defined by duality; if $u\in\CmI(X;E)$ and $\kappa:W\to\kappa(W)$ is a
trivializing chart, then for $\psi\in\CIc(W)$ the localization $\psi u$ is
represented by an $N$-tuple of compactly supported distributions on $\kappa(W)$,
which we write $\kappa_*(\psi u)=(\kappa^{-1})^*(\psi u)$.

Distributional sections can equally be defined ``algebraically''.

\begin{proposition}\label{manifolds:prop:tensor}
For any vector bundle $U$ over the compact manifold $Y$ there is a natural
isomorphism
\begin{equation*}
\CmI(Y)\otimes_{\CI(Y)}\CI(Y;U)\cong\CmI(Y;U).
\end{equation*}
\end{proposition}
\begin{proof}
The map $\sum_iu_i\otimes g_i\mapsto\sum_i u_ig_i$, where
$(u g)(\psi)=u(g\cdot\psi)$ for $\psi\in\CI(Y;U'\otimes\Omega)$ (so that
$g\cdot\psi\in\CI(Y;\Omega)$), is well defined on the tensor product since it is
$\CI(Y)$-balanced. To invert it, choose finitely many sections $g_i\in\CI(Y;U)$
and $q_i\in\CI(Y;U')$ with
\begin{equation}\label{manifolds:eq:frame-identity}
\sum_i g_i\otimes q_i=\Id_U\in\CI(Y;\End(U)).
\end{equation}
These exist: take a finite partition of unity $\rho_k$ subordinate to a cover by
trivializing open sets $U_k$, local frames $e_{k,l}$ over $U_k$ with dual frames
$e^{k,l}$, functions $\psi_k\in\CIc(U_k)$ equal to one on $\supp\rho_k$, and use
the sections $\rho_ke_{k,l}$ and $\psi_ke^{k,l}$, extended by zero. Then
$\sum_k\sum_l\rho_ke_{k,l}\otimes\psi_ke^{k,l}=\sum_k\rho_k\Id=\Id$. Define
$u\mapsto\sum_i(u\cdot q_i)\otimes g_i$ for $u\in\CmI(Y;U)$, where $u\cdot q_i\in\CmI(Y)$
is the pairing. By \eqref{manifolds:eq:frame-identity},
$\sum_i(u\cdot q_i)g_i=u$, while for $u=\sum_ju_j\otimes h_j$,
$\sum_i\sum_ju_j(h_j\cdot q_i)\otimes g_i=\sum_ju_j\otimes\sum_i(h_j\cdot q_i)g_i
=\sum_ju_j\otimes h_j$, using the $\CI(Y)$-balancing. So the two maps are
mutually inverse.
\end{proof}

The same argument shows that any $\CI(Y)$-module of ``distributions'' which is
local and invariant under multiplication by smooth functions, such as the
conormal spaces below, can be extended to sections of bundles by tensoring.

\section{Euclidean and conormal preliminaries}\label{manifolds:sec:recall}

\subsection{The calculus on $\R^n$}
We collect the results from \cref{ch:euclidean,ch:microlocal} which are used
below. Here $\Psido{m}(\R^n)$ is the space of operators $\Op(a)$,
$a\in S^m(\R^n;\R^n)$, with kernels
$K(x,y)=\int e^{i(x-y)\cdot\xi}a(x,\xi)\,\dbar\xi$.
\begin{enumerate}[label=(E\arabic*),leftmargin=3em]
\item\label{manifolds:E1} $\Psido{m}(\R^n)\subset\Psido{m'}(\R^n)$ for $m\le m'$;
$\Psido{-\infty}(\R^n)=\bigcap_m\Psido m(\R^n)$ consists of operators with smooth
kernels, and every operator with kernel in $\CIc(\R^{2n})$ lies in
$\Psido{-\infty}(\R^n)$. Kernels of elements of $\Psido m(\R^n)$ are smooth off the
diagonal.
\item\label{manifolds:E2} (Amplitudes) If $a(x,y,\xi)$ satisfies
$|D^\alpha_xD^\beta_yD^\gamma_\xi a|\le C\ang\xi^{m-|\gamma|}$ then the operator with
kernel $\int e^{i(x-y)\cdot\xi}a(x,y,\xi)\,\dbar\xi$ lies in $\Psido m(\R^n)$, with
principal symbol $a(x,x,\xi)$. Equivalently, every $A\in\Psido m(\R^n)$ has a
\emph{left} symbol $a_L(x,\xi)$ and a \emph{right} symbol $a_R(y,\xi)$, both in
$S^m$, with $\sigma_m(A)=[a_L]=[a_R]$. In particular $\phi A\psi\in\Psido m(\R^n)$ for
$\phi,\psi\in\CIb(\R^n)$, and the kernel of $A$ may be multiplied by $\chi(x-y)$,
$\chi\in\CIc(\R^n)$ equal to one near $0$, modulo $\Psido{-\infty}(\R^n)$.
\item\label{manifolds:E3} (Composition and adjoints) $\Psido m(\R^n)\circ\Psido{m'}(\R^n)\subset\Psido{m+m'}(\R^n)$ with
$\sigma_{m+m'}(AB)=\sigma_m(A)\sigma_{m'}(B)$; the $L^2$ adjoint satisfies
$A^*\in\Psido m(\R^n)$, $\sigma_m(A^*)=\overline{\sigma_m(A)}$.
\item\label{manifolds:E4} (Principal symbol) The sequence
$0\to\Psido{m-1}(\R^n)\to\Psido m(\R^n)\xrightarrow{\sigma_m}S^m/S^{m-1}\to0$ is exact,
with $\sigma_m(\Op(a))=[a]$.
\item\label{manifolds:E5} (Asymptotic completeness) If $A_j\in\Psido{m-j}(\R^n)$
there exists $A\in\Psido m(\R^n)$ with $A-\sum_{j<N}A_j\in\Psido{m-N}(\R^n)$ for all
$N$. If the kernels of the $A_j$ have support in a fixed compact set $K$, $A$ may
be chosen with kernel supported in any given neighbourhood of $K$ (replace $A$ by
$\chi A\chi$ with $\chi\in\CIc(\R^n)$ equal to one near both projections of $K$).
\item\label{manifolds:E6} (Sobolev boundedness) Each $A\in\Psido m(\R^n)$ is bounded
$\Hs s(\R^n)\to\Hs{s-m}(\R^n)$ for every $s\in\R$; in particular $\Psido0(\R^n)$
acts boundedly on $L^2(\R^n)$, and multiplication by $\phi\in\CIb(\R^n)$ is bounded on
every $\Hs s(\R^n)$.
\item\label{manifolds:E7} (Coordinate invariance) Let $\Phi:\Omega\to\Omega'$ be a
diffeomorphism between open subsets of $\R^n$ and $A\in\Psido m(\R^n)$ have kernel
with compact support in $\Omega\times\Omega$. Then
$\Phi_*A=(\Phi^{-1})^*\circ A\circ\Phi^*$, extended by zero, lies in
$\Psido m(\R^n)$, and
\begin{equation}\label{manifolds:eq:symbol-invariance}
\sigma_m(\Phi_*A)(\Phi(x),\eta)=\sigma_m(A)\big(x,(d\Phi_x)^T\eta\big).
\end{equation}
\end{enumerate}
The symbol formula \eqref{manifolds:eq:symbol-invariance} follows from the
proof of coordinate invariance: writing $G=\Phi^{-1}$ and
$G(x)-G(y)=T(x,y)(x-y)$, the kernel of $\Phi_*A$ near the diagonal is given by the
amplitude $c(x,y,\eta)=a(G(x),T(x,y)^{-T}\eta)\,|\det dG(y)|\,|\det T(x,y)|^{-1}$, and
since $T(x,x)=dG(x)$ its value on the diagonal is $a(G(x),(dG_x)^{-T}\eta)$, which
is \eqref{manifolds:eq:symbol-invariance} because $(dG_{\Phi(x')})^{-T}=(d\Phi_{x'})^T$.
Formula \eqref{manifolds:eq:symbol-invariance} says precisely that the principal
symbol is invariantly a function on the cotangent bundle: it is the pull-back
under the induced map $T^*\Omega'\to T^*\Omega$.

\subsection{Conormal distributions}\label{manifolds:sec:conormal}
We recall the definition from \cref{ch:conormal}. Let $M$ be a manifold of
dimension $N$ and $Y\subset M$ a closed embedded submanifold of codimension $k$.
The space $\Ic m(M,Y)$ of distributions (or distributional densities) conormal
to $Y$ of order $m$ consists of those $u$ which are smooth on $M\setminus Y$ and
such that, near each point of $Y$, in local coordinates $(w,z)\in\R^{N-k}\times\R^k$
in which $Y=\{z=0\}$,
\begin{equation}\label{manifolds:eq:conormal-local}
u=\int_{\R^k}e^{iz\cdot\zeta}a(w,\zeta)\,\dbar\zeta+v,\quad
a\in S^{m+\frac N4-\frac k2}(\R^{N-k};\R^k),\ v\ \text{smooth}.
\end{equation}
(For distributional densities the formula is multiplied by the coordinate
density $|dw\,dz|$.) The subspace in which the amplitudes $a$ are classical
symbols is denoted $\Ic m_{\mathrm{ph}}(M,Y)$. We use the following properties.
\begin{enumerate}[label=(C\arabic*),leftmargin=3em]
\item\label{manifolds:C1} $\Ic m(M,Y)$ is a $\CI(M)$-module, is local, and is
invariant under diffeomorphisms: if $\Phi:M\to M'$ is a diffeomorphism with
$\Phi(Y)=Y'$ then $\Phi^*\Ic m(M',Y')=\Ic m(M,Y)$. Moreover
$\Ic m\subset\Ic{m'}$ for $m\le m'$ and $\bigcap_m\Ic m(M,Y)=\CI(M)$.
\item\label{manifolds:C2} The amplitude $a$ in \eqref{manifolds:eq:conormal-local}
may be allowed to depend on $z$ as well, with symbol estimates uniform on compact
sets in $(w,z)$, without changing the space.
\item\label{manifolds:C3} Differential operators of order $p$ map $\Ic m(M,Y)$ into
$\Ic{m+p}(M,Y)$, and multiplication by a smooth function vanishing on $Y$ maps
$\Ic m(M,Y)$ into $\Ic{m-1}(M,Y)$.
\end{enumerate}
The order normalization is such that the Dirac delta at a point of $\R^N$ has
order $N/4$. For a vector bundle $E$ over $M$ we set
\begin{equation*}
\Ic m(M,Y;E)=\Ic m(M,Y)\otimes_{\CI(M)}\CI(M;E),
\end{equation*}
the space of finite sums $\sum_iu_ie_i$, $u_i\in\Ic m(M,Y)$, $e_i\in\CI(M;E)$, modulo
the relations generated by $\CI(M)$-linearity. As in
\cref{manifolds:prop:tensor}, if $u$ has support in an open set over which $E$
is trivial, with local frame $e_i$, then $u=\sum_iu_ie_i$ with
$u_i\in\Ic m(M,Y)$, and there are natural inclusions
\begin{equation*}
\CI(M;E)\hookrightarrow\Ic m(M,Y;E)\hookrightarrow\CmI(M;E).
\end{equation*}

\subsubsection*{Integration, restriction and push-forward}
If $M$ is compact and $V$ is a fixed vector space, the integral of a
distributional density, i.e.\ its pairing with the constant function $1$, gives
a linear map $\int_M:\Ic m(M,Y;V\otimes\Omega_M)\to V$.

\begin{proposition}\label{manifolds:prop:restriction}
Let $Z\subset M$ be an embedded submanifold transversal to $Y$, i.e.\
$T_pY+T_pZ=T_pM$ for all $p\in Y\cap Z$. Then restriction
$\CI(M;E)\to\CI(Z;E|_Z)$ extends to a map
\begin{equation*}
\Ic m(M,Y;E)\longrightarrow\Ic{m+\frac14\operatorname{codim}Z}(Z,Y\cap Z;E|_Z).
\end{equation*}
\end{proposition}
\begin{proof}
Elements supported away from $Y$ are smooth, so by a partition of unity it
suffices to consider $u$ supported near a point $p\in Y\cap Z$, where $E$ may be
taken trivial. Let $z_1,\dots,z_k$ be local defining functions for $Y$ and
$t_1,\dots,t_q$ for $Z$, $q=\operatorname{codim}Z$. Transversality means
precisely that the differentials of the $z_j$ and $t_l$ are linearly independent at
$p$, so they can be completed by functions $w'$ to a coordinate system $(w',t,z)$
near $p$, in which $Y=\{z=0\}$ and $Z=\{t=0\}$. In these coordinates
$u=\int e^{iz\cdot\zeta}a(w',t,\zeta)\dbar\zeta+v$ with
$a\in S^{m+\frac N4-\frac k2}$, and $u$ depends smoothly on $t$ as a distribution
in $(w',z)$. We set
$u|_Z=\int e^{iz\cdot\zeta}a(w',0,\zeta)\dbar\zeta+v|_{t=0}$; this agrees with the
restriction if $a\in S^{-\infty}$, and in general is the limit of the restrictions of
the smooth functions obtained by replacing $a$ by $a(w',t,\zeta)\chi(\zeta/R)$,
$\chi\in\CIc(\R^k)$ equal to one near $0$, as $R\to\infty$; hence it is independent
of choices. Since $\dim Z=N-q$, the order $m''$ of $u|_Z$ is fixed by
$m+\frac N4-\frac k2=m''+\frac{N-q}4-\frac k2$, i.e.\ $m''=m+\frac q4$.
\end{proof}

Now let $\phi:M\to B$ be a fibration, locally the projection of a product,
transversal to $Y$ in the sense that each fibre is transversal to $Y$:
$T_pY+T_p\phi^{-1}(\phi(p))=T_pM$ for $p\in Y$. Let $\Omega_\phi$ be the bundle
of densities on the fibres, so that $\Omega_M=\Omega_\phi\otimes\phi^*\Omega_B$.

\begin{proposition}\label{manifolds:prop:pushforward}
If $\phi:M\to B$ is a fibration with compact fibres transversal to $Y$, and $E$
is a vector bundle over $B$, then fibre integration extends from smooth sections
to a linear map
\begin{equation*}
\phi_*:\Ic m(M,Y;\phi^*E\otimes\Omega_\phi)\longrightarrow\CI(B;E).
\end{equation*}
\end{proposition}
\begin{proof}
For smooth sections, locally in $B$ where $M=U\times Z$, $\phi_*u(b)=\int_Zu(b,z)$,
which is smooth in $b$. By a partition of unity it suffices to consider $u$
supported near a point $p\in Y$ and $E$ trivial. By transversality, the defining
functions of $Y$ restrict to each fibre near $p$ to functions with independent
differentials, so we may choose coordinates $(b,z',z'')$ near $p$ in which $\phi$
is $(b,z',z'')\mapsto b$ and $Y=\{z'=0\}$. Then, modulo a smooth term,
$u=\int e^{iz'\cdot\zeta}a(b,z'',\zeta)\dbar\zeta\,|dz'\,dz''|$ with compact support in
$(b,z'')$. Choosing $\chi\in\CIc$ equal to one on the support of $u$ in $z'$,
\begin{equation*}
\phi_*u(b)=\int\!\!\int\chi(z')\int e^{iz'\cdot\zeta}a(b,z'',\zeta)\,\dbar\zeta\,dz'\,dz''
=\int\!\!\int a(b,z'',\zeta)\hat\chi(-\zeta)\,\dbar\zeta\,dz'',
\end{equation*}
which is smooth in $b$ since $\hat\chi\in\cS$ and $a$ is a symbol. (Formally,
$\int e^{iz'\cdot\zeta}dz'=(2\pi)^k\delta(\zeta)$ and $\phi_*u(b)=\int a(b,z'',0)\,dz''$.)
\end{proof}

\section{Pseudodifferential operators on a compact manifold}\label{manifolds:sec:definition}

\subsection{Definition}
Let $E$ and $F$ be complex vector bundles over the compact manifold $X$. The
diagonal $\Diag=\{(x,x)\}\subset X^2$ is an embedded submanifold of codimension
$n$, and the fibres $\{x\}\times X$ of $\pi_L$ are transversal to it.

\begin{definition}\label{manifolds:def:psido}
For $m\in\R$, the space of pseudodifferential operators of order $m$ from
sections of $E$ to sections of $F$ is
\begin{equation}\label{manifolds:eq:def-psido}
\Psido m(X;E,F)=\Ic m(X^2,\Diag;\Hom(E,F)\otimes\Omega_R),
\end{equation}
with $\Hom(E,F)$ as in \eqref{manifolds:eq:big-hom}. The operator defined by
$A\in\Psido m(X;E,F)$ is
\begin{equation}\label{manifolds:eq:action}
Au(x)=\int_XA(x,y)u(y),\quad\text{i.e.}\quad
Au=(\pi_L)_*\big(A\cdot\pi_R^*u\big),\qquad u\in\CI(X;E).
\end{equation}
We write $\Psido m(X;E)=\Psido m(X;E,E)$, $\Psido m(X)$ for trivial line bundles,
and $\Psido m_{\mathrm{ph}}$ for the classical operators, defined using
$\Ic m_{\mathrm{ph}}$.
\end{definition}

In \eqref{manifolds:eq:action}, $\pi_R^*u\in\CI(X^2;\pi_R^*E)$, the product with
the kernel contracts $E'$ with $E$ and so gives an element of
$\Ic m(X^2,\Diag;\pi_L^*F\otimes\Omega_R)$; since $\Omega_R$ is the density bundle
of the fibres of $\pi_L$, which are transversal to $\Diag$,
\cref{manifolds:prop:pushforward} shows that $Au\in\CI(X;F)$. Thus
$A:\CI(X;E)\to\CI(X;F)$, and the operator determines its kernel, by the Schwartz
kernel theorem; we identify the two, and also write $K_A$ for the kernel. For
$\phi,\psi\in\CI(X)$, the kernel of $\phi A\psi$ is $\phi(x)A(x,y)\psi(y)$.

The space of \emph{smoothing operators} is
\begin{equation}\label{manifolds:eq:smoothing}
\Psido{-\infty}(X;E,F)=\bigcap_m\Psido m(X;E,F)=\CI(X^2;\Hom(E,F)\otimes\Omega_R),
\end{equation}
by \ref{manifolds:C1}. Such an operator, $Au(x)=\int_XA(x,y)u(y)$, extends to a
continuous map $\CmI(X;E)\to\CI(X;F)$, by letting the distribution $u$ act in
the variable $y$.

\subsection{Local description}
We now connect \cref{manifolds:def:psido} with the calculus on $\R^n$. If
$\kappa:W\to\kappa(W)$ is a chart over which $E$ and $F$ are trivialized, and $A$
has kernel with compact support in $W\times W$, its \emph{local representative}
is the matrix of operators
\begin{equation*}
\kappa_*A=(\kappa^{-1})^*\circ A\circ\kappa^*:\CIc(\kappa(W))^{N_E}\longrightarrow\CIc(\kappa(W))^{N_F},
\end{equation*}
where the right density factor is trivialized by the coordinate density.

\begin{theorem}\label{manifolds:thm:local}
A linear operator $A:\CI(X;E)\to\CI(X;F)$ is given by an element of
$\Psido m(X;E,F)$ if and only if
\begin{enumerate}[label=\textup{(\roman*)}]
\item if $\phi,\psi\in\CI(X)$ have disjoint supports then $\phi A\psi$ has
kernel in $\CI(X^2;\Hom(E,F)\otimes\Omega_R)$, and
\item for every trivializing chart $\kappa:W\to\kappa(W)$ and every
$\psi\in\CIc(W)$, the entries of $\kappa_*(\psi A\psi)$ are elements of
$\Psido m(\R^n)$ (with kernels supported in $\kappa(\supp\psi)^2$).
\end{enumerate}
Moreover it suffices to require \textup{(ii)} for the charts of one finite cover
of $X$ by trivializing charts.
\end{theorem}
\begin{proof}
Suppose $A\in\Psido m(X;E,F)$. Its kernel is smooth outside the diagonal, which
gives (i). For (ii), $\psi(x)\psi(y)A(x,y)$ is conormal of order $m$ with compact
support in $W\times W$, so by \ref{manifolds:C1} the kernel of $\kappa_*(\psi A\psi)$
is a matrix of compactly supported elements of $\Ic m(\R^{2n},\Diag)$. Use the
linear coordinates $(x,z)=(x,x-y)$ on $\R^{2n}$, in which $\Diag=\{z=0\}$; here
$N=2n$, $k=n$ and the symbol order in \eqref{manifolds:eq:conormal-local} is
$m+\frac{2n}4-\frac n2=m$. By \ref{manifolds:C1} the local representation
\eqref{manifolds:eq:conormal-local} holds in these coordinates, so with a partition
of unity in $x$ and a cutoff $\chi(z)$ equal to one near $z=0$, each entry of the
kernel is a finite sum of terms
$\chi_j(x)\chi(x-y)\int e^{i(x-y)\cdot\xi}a_j(x,\xi)\dbar\xi$, $a_j\in S^m$, plus a
function in $\CIc(\R^{2n})$. By \ref{manifolds:E2} and \ref{manifolds:E1} these
are kernels of elements of $\Psido m(\R^n)$.

Conversely suppose (i) holds and (ii) holds for the charts $\kappa_i:W_i\to\kappa_i(W_i)$ of
a finite cover. By (i), near any point $(\bar x,\bar y)$ with $\bar x\ne\bar y$ the
kernel of $A$ coincides with that of $\phi A\psi$, where $\phi=1$ near $\bar x$,
$\psi=1$ near $\bar y$ and $\supp\phi\cap\supp\psi=\emptyset$; so $K_A$ is smooth off
the diagonal. Near $(\bar x,\bar x)$, with $\bar x\in W_i$ and $\psi\in\CIc(W_i)$ equal
to one near $\bar x$, $K_A$ coincides with the kernel of $\psi A\psi$, whose local
representative is $\int e^{i(x-y)\cdot\xi}a(x,\xi)\dbar\xi$ with $a\in S^m$: this is
of the form \eqref{manifolds:eq:conormal-local} in the coordinates $(x,x-y)$. Since
conormality is a local property, $K_A\in\Ic m(X^2,\Diag;\Hom(E,F)\otimes\Omega_R)$.
In particular (ii) then holds for all trivializing charts.
\end{proof}

The second half of the proof, combined with the first, shows again that the
compactly supported elements of $\Psido m(\R^n)$ are invariant under
diffeomorphisms, as in \ref{manifolds:E7}. It also gives the following, which
we use repeatedly to build operators on $X$ out of operators on $\R^n$.

\begin{corollary}\label{manifolds:cor:local-ops}
If $\kappa:W\to\kappa(W)$ is a trivializing chart and $B$ is an
$N_F\times N_E$ matrix of elements of $\Psido m(\R^n)$ with kernels of compact
support in $\kappa(W)\times\kappa(W)$, then $\kappa^*\circ B\circ(\kappa^{-1})^*$,
extended by zero, is an element of $\Psido m(X;E,F)$.
\end{corollary}

\subsection{Normal fibrations and reduced forms}
To write the kernel of an operator supported in a coordinate patch as an
inverse Fourier transform, one identifies a neighbourhood of the diagonal with
its normal bundle. The normal bundle of $\Diag\subset\R^{2n}$ is identified with
$T\R^n=\R^n\times\R^n$ via the left tangent vectors $(w,0)$; there remain
several natural choices of a \emph{normal fibration}, i.e.\ a linear map
$\R^{2n}\to T\R^n$ inducing this identification. The \emph{left}, \emph{right}
and \emph{Weyl} fibrations are
\begin{equation*}
f_L(x,y)=(x,x-y),\quad f_R(x,y)=(y,x-y),\quad f_W(x,y)=\big(\tfrac{x+y}2,x-y\big).
\end{equation*}
In each case the fibre variable is $x-y$ and one of $x$, $y$ or $\frac{x+y}2$ is held
fixed. Correspondingly an operator $A\in\Psido m(\R^n)$ with compactly supported
kernel has the three representations
\begin{equation}\label{manifolds:eq:three-reps}
K_A=\int e^{i(x-y)\cdot\xi}a_L(x,\xi)\,\dbar\xi
=\int e^{i(x-y)\cdot\xi}a_R(y,\xi)\,\dbar\xi
=\int e^{i(x-y)\cdot\xi}a_W\big(\tfrac{x+y}2,\xi\big)\,\dbar\xi,
\end{equation}
with $a_L,a_R,a_W\in S^m$ (classical if $A$ is): the first two exist by
\ref{manifolds:E2}, and the Weyl symbol is obtained in the same way by Taylor
expansion of an amplitude about $x=y$ (see \cite[Section~18.5]{Hormander-III}).
All three have the same principal symbol, and they may be taken to have compact
support in the base variable.

\begin{lemma}\label{manifolds:lem:reduced}
If $A\in\Psido m(\R^n)$ has compactly supported kernel, then for $f\in\cS(\R^n)$
\begin{equation}\label{manifolds:eq:left-right}
Af(x)=\int e^{ix\cdot\xi}a_L(x,\xi)\hat f(\xi)\,\dbar\xi,\qquad
\widehat{Af}(\xi)=\int e^{-iy\cdot\xi}a_R(y,\xi)f(y)\,dy.
\end{equation}
\end{lemma}
\begin{proof}
The first formula is the definition of $\Op(a_L)$ after inserting
$\hat f(\xi)=\int e^{-iy\cdot\xi}f(y)dy$. For the second, let $g(\xi)$ denote the
right side, with $a_R$ compactly supported in $y$; integration by parts in $y$
shows that $|\xi^\alpha g(\xi)|\le C_\alpha\ang\xi^m$, so $g$ is smooth and
polynomially bounded, and by Fubini (first for $a_R$ replaced by
$a_R(y,\xi)\chi(\xi/R)$, then letting $R\to\infty$)
$Af(x)=\int e^{ix\cdot\xi}g(\xi)\,\dbar\xi$, i.e.\ $\widehat{Af}=g$.
\end{proof}

\subsection{Differential operators}
A differential operator $P\in\Diff k(X;E,F)$ is a linear map
$\CI(X;E)\to\CI(X;F)$ which in any local coordinates and trivializations takes the
form
\begin{equation}\label{manifolds:eq:diffop}
P=\sum_{|\alpha|\le k}p_\alpha(x)D_x^\alpha,\quad p_\alpha \text{ smooth }
N_F\times N_E \text{ matrices}.
\end{equation}
By \cref{manifolds:thm:local}, $\Diff k(X;E,F)\subset\Psido k_{\mathrm{ph}}(X;E,F)$:
condition (i) holds trivially since $\phi P\psi=0$ if $\supp\phi\cap\supp\psi=\emptyset$,
and $\psi P\psi=\sum\psi p_\alpha D^\alpha\psi$ in local coordinates is a composite of
elements of $\Psido{}(\R^n)$. The kernel of $P$ is
$\sum_\alpha p_\alpha(x)D^\alpha_x\delta(x-y)|dy|$, conormal to the diagonal.

The leading part of \eqref{manifolds:eq:diffop} has an invariant meaning. Let
$\bar x\in X$, $\Xi\in T^*_{\bar x}X$, choose a real-valued $f\in\CI(X)$ with
$df(\bar x)=\Xi$ and $u\in\CI(X;E)$. Then by Leibniz' formula
\begin{equation}\label{manifolds:eq:oscillatory-test}
e^{-itf}P(e^{itf}u)=\sum_{j=0}^kt^jU_j,\qquad
U_k(\bar x)=\sum_{|\alpha|=k}p_\alpha(\bar x)\,\Xi^\alpha\,u(\bar x),
\end{equation}
where $\Xi^\alpha=(\pa f(\bar x))^\alpha$ in local coordinates; the left side is
defined invariantly, so the coefficient of $t^k$ is too. Hence the
\emph{principal symbol}
\begin{equation*}
\sigma_k(P)(x,\xi)=\sum_{|\alpha|=k}p_\alpha(x)\xi^\alpha\in\CI(T^*X;\hom(E,F))
\end{equation*}
is a well-defined section, homogeneous polynomial of degree $k$ on the fibres;
$P$ is called \emph{elliptic} if $\sigma_k(P)(x,\xi)$ is invertible for all
$\xi\ne0$. For example, for the exterior derivative
$d(e^{itf}u)=e^{itf}(it\,df\wedge u+du)$, so $\sigma_1(d)(x,\xi)=i\,\xi\wedge$.

\section{The calculus}\label{manifolds:sec:calculus}

\subsection{Localization}
\begin{lemma}\label{manifolds:lem:localization}
Let $\{U_i\}$ be a finite open cover of $X$ and $A\in\Psido m(X;E,F)$. Then
\begin{equation}\label{manifolds:eq:localization}
A=\sum_iA_i+A',\qquad A_i\in\Psido m(X;E,F),\ \supp K_{A_i}\Subset U_i\times U_i,\
A'\in\Psido{-\infty}(X;E,F).
\end{equation}
\end{lemma}
\begin{proof}
Let $\{\phi_a\}$ be a partition of unity as in \cref{manifolds:lem:good-partition}.
Then $A=\sum_{a,b}\phi_aA\phi_b$, the kernel of $\phi_aA\phi_b$ being
$\phi_a(x)A(x,y)\phi_b(y)$. If $\supp\phi_a\cap\supp\phi_b=\emptyset$ this kernel
vanishes near the diagonal, so is smooth. Otherwise
$\supp\phi_a\cup\supp\phi_b\subset U_i$ for some $i$, and the term has kernel
supported in $U_i\times U_i$. Collecting terms gives \eqref{manifolds:eq:localization}.
\end{proof}

Taking the $U_i$ to be trivializing charts, each $A_i$ in
\eqref{manifolds:eq:localization} is given by a matrix of compactly supported
elements of $\Psido m(\R^n)$, and conversely by \cref{manifolds:cor:local-ops}.

\begin{proposition}\label{manifolds:prop:continuity}
Each $A\in\Psido m(X;E,F)$ is continuous $\CI(X;E)\to\CI(X;F)$. If
$R\in\Psido{-\infty}(X;G,E)$ then $AR\in\Psido{-\infty}(X;G,F)$.
\end{proposition}
\begin{proof}
In \eqref{manifolds:eq:localization}, $A'$ has a smooth kernel, so is continuous,
and each $A_i$ is, in local coordinates, a matrix of elements of $\Psido m(\R^n)$
with compactly supported kernels; these are continuous on $\cS(\R^n)$ and hence on
smooth functions supported in a fixed compact set. The kernel of $AR$ is
$(x,z)\mapsto A(K_R(\cdot,z))(x)$; since $z\mapsto K_R(\cdot,z)$ is a smooth map into
$\CI(X;E)$ (after trivializing the bundles in $z$ locally) and $A$ is continuous and
linear, this is smooth in $(x,z)$.
\end{proof}

\subsection{Adjoints, transposes and action on distributions}
Fix a smooth positive density $\nu$ on $X$ and hermitian metrics on $E$ and $F$,
so that $L^2(X;E)$ and $L^2(X;F)$ are defined by \eqref{manifolds:eq:L2}.

\begin{proposition}\label{manifolds:prop:adjoint}
If $A\in\Psido m(X;E,F)$ there is a unique $A^*\in\Psido m(X;F,E)$ such that
$\inner{Au}v=\inner u{A^*v}$ for all $u\in\CI(X;E)$, $v\in\CI(X;F)$, and
$\sigma_m(A^*)=\sigma_m(A)^*$ (pointwise adjoint with respect to the metrics).
\end{proposition}
\begin{proof}
Write $A(x,y)=k(x,y)\nu(y)$ and set $A^*(x,y)=k(y,x)^*\nu(y)$, where $k(y,x)^*$ is
the adjoint of $k(y,x)\in\hom(E_x,F_y)$. The reflection $(x,y)\mapsto(y,x)$
preserves the diagonal, and $\Ic m$ is invariant under complex conjugation
(replace $a(w,\zeta)$ by $\overline{a(w,-\zeta)}$ in \eqref{manifolds:eq:conormal-local}),
so by \ref{manifolds:C1} $A^*\in\Psido m(X;F,E)$. The identity
$\inner{Au}v=\inner u{A^*v}$ is the statement that both sides equal the pairing of
the kernel with $\overline{v}(x)\otimes u(y)$. Uniqueness is clear. In local
coordinates where $\nu=\nu_0(x)|dx|$ and the metrics are given by positive matrices
$h_E,h_F$, the local representative of $A^*$ is $\nu_0^{-1}h_E^{-1}B^*h_F\nu_0$ where
$B^*$ is the Euclidean adjoint of that of $A$, so \ref{manifolds:E3} gives
$\sigma_m(A^*)=h_E^{-1}\sigma_m(A)^{\dagger}h_F$, which is the adjoint with respect to the
metrics.
\end{proof}

\begin{corollary}\label{manifolds:cor:smoothing-ideal}
If $A\in\Psido m(X;F,G)$ and $R\in\Psido{-\infty}(X;E,F)$ then $AR$ and $R'A$ are
smoothing, for any $R'\in\Psido{-\infty}(X;G,E)$.
\end{corollary}
\begin{proof}
$AR$ is smoothing by \cref{manifolds:prop:continuity}, and $R'A=(A^*R'^*)^*$, the
adjoint of a smoothing operator being smoothing.
\end{proof}

Similarly the \emph{transpose} of $A$, with kernel $A^t(x,y)=A(y,x)^t$, is an
element of $\Psido m(X;F'\otimes\Omega,E'\otimes\Omega)$ (the density factors
match, since $A(y,x)\in F_y\otimes E'_x\otimes\Omega_x$), and
$\int_X(Au)\cdot\phi=\int_Xu\cdot A^t\phi$ for smooth $u$ and $\phi$. We may therefore
extend $A$ to distributions by
\begin{equation}\label{manifolds:eq:distributions-action}
A:\CmI(X;E)\longrightarrow\CmI(X;F),\qquad (Au)(\phi)=u(A^t\phi),
\end{equation}
which is continuous and, by density, the unique continuous extension of $A$ on
$\CI(X;E)$. All the identities below between operators on smooth sections
therefore hold on distributions.

\subsection{Composition}
\begin{theorem}\label{manifolds:thm:composition}
If $E$, $F$, $G$ are vector bundles over the compact manifold $X$ then
\begin{equation*}
\Psido m(X;F,G)\circ\Psido{m'}(X;E,F)\subset\Psido{m+m'}(X;E,G).
\end{equation*}
\end{theorem}
\begin{proof}
Let $A\in\Psido m(X;F,G)$, $B\in\Psido{m'}(X;E,F)$. Choose a finite cover by charts
$U_c$ over which $E$, $F$ and $G$ are trivial, a subordinate partition of unity
$\phi_c$, and $\chi_c,\chi'_c,\tilde\chi_c\in\CIc(U_c)$ with $\chi_c=\chi'_c=1$ near
$\supp\phi_c$ and $\tilde\chi_c=1$ near $\supp\chi_c\cup\supp\chi'_c$. Then
\begin{equation*}
AB=\sum_cA\phi_cB=\sum_c\big(\chi_cA\phi_cB\chi'_c+(1-\chi_c)A\phi_cB+\chi_cA\phi_cB(1-\chi'_c)\big).
\end{equation*}
Since $(1-\chi_c)$ and $\phi_c$ have disjoint supports, $(1-\chi_c)A\phi_c$ is smoothing,
and so is its product with $B$ by \cref{manifolds:cor:smoothing-ideal}; similarly for
the last term. Finally
$\chi_cA\phi_cB\chi'_c=\chi_c(\tilde\chi_cA\tilde\chi_c)\phi_c(\tilde\chi_cB\tilde\chi_c)\chi'_c$
is, in the chart $U_c$, a product of matrices of elements of $\Psido m(\R^n)$ and
$\Psido{m'}(\R^n)$ with compactly supported kernels, so lies in
$\Psido{m+m'}(\R^n)$ by \ref{manifolds:E3}, and the result follows from
\cref{manifolds:cor:local-ops}.
\end{proof}

The local step can be seen directly from \cref{manifolds:lem:reduced}: writing
the first factor in left-reduced and the second in right-reduced form, for
$f\in\cS(\R^n)$,
\begin{equation}\label{manifolds:eq:left-right-composition}
\begin{gathered}
ABf(x)=\int e^{ix\cdot\xi}a_L(x,\xi)\Big(\int e^{-iy\cdot\xi}b_R(y,\xi)f(y)\,dy\Big)\dbar\xi,\\
K_{AB}=\int e^{i(x-y)\cdot\xi}a_L(x,\xi)b_R(y,\xi)\,\dbar\xi .
\end{gathered}
\end{equation}
The amplitude $a_L(x,\xi)b_R(y,\xi)$ satisfies the estimates of
\ref{manifolds:E2} with order $m+m'$, so $AB\in\Psido{m+m'}(\R^n)$ with principal
symbol $a_L(x,\xi)b_R(x,\xi)$. Equivalently, $K_{AB}$ is the restriction to $z=y$ of
$\int e^{i(x-y)\cdot\xi}a_L(x,\xi)b_R(z,\xi)\dbar\xi$, which is conormal in
$\R^{3n}$ to $\{x=y\}$; the submanifold $\{z=y\}$ is transversal to it, and
\cref{manifolds:prop:restriction} shows again that the result is conormal.

\begin{remark}\label{manifolds:rem:global-composition}
When one factor is smoothing, the composite can also be treated globally. With
$\pi_S(x,y,z)=(x,y)$, $\pi_F(x,y,z)=(y,z)$ and $\pi_C(x,y,z)=(x,z)$ on $X^3$,
the kernel of $AR$, $R$ smoothing, is
$(\pi_C)_*(\pi_S^*A\cdot\pi_F^*R)$, where the density factor of $A$ in $y$ is
used for the integration. The pull-back $\pi_S^*A$ is conormal to
$\pi_S^{-1}(\Diag)=\{x=y\}$, of order $m-\frac n4$ by the behaviour of conormal
distributions under pull-back by submersions (\cref{ch:conormal}), the product with
the smooth section $\pi_F^*R$ preserves this, and the fibres of $\pi_C$ (on
which $y$ varies) are transversal to $\{x=y\}$, so the push-forward is smooth
by \cref{manifolds:prop:pushforward}. For general $A$ and $B$, however, the
corresponding product $\pi_S^*A\cdot\pi_F^*B$ is not defined in this simple
way, since both factors are singular at the triple diagonal; this is why the
proof above localizes.
\end{remark}

\subsection{Symbols and the principal symbol}
Let $V\to X$ be a real vector bundle (we only need $V=T^*X$) and $W$ a complex
vector bundle over $X$, pulled back to $V$. The space $S^m(V;W)$ consists of
the smooth sections $a$ of $W$ over $V$ such that, in any local trivializations of
$V$ and $W$ over a coordinate patch, and on any compact subset $K$ of the patch,
\begin{equation*}
|D^\alpha_xD^\beta_\xi a(x,\xi)|\le C_{K,\alpha,\beta}\ang\xi^{m-|\beta|}.
\end{equation*}
This is independent of the trivializations: under
$(x,\xi)\mapsto(x,L(x)\xi)$, with $L$ smooth and invertible, and a change
$a\mapsto g(x)a$ of trivialization of $W$, the chain rule shows that each
$\xi$-derivative of the new expression is a combination of $\xi$-derivatives of
the old one with smooth coefficients, and each $x$-derivative produces terms
$(\pa_xL)\xi\cdot\pa_\xi a$ which are of the same order; since $\ang{L\xi}$ and
$\ang\xi$ are comparable on compact sets the estimates follow by induction. Since $X$
is compact, $S^m(V;W)$ is a Fr\'echet space. Set
$S^{m-[1]}(V;W)=S^m(V;W)/S^{m-1}(V;W)$. The classical symbols $S^m_{\mathrm{ph}}$
are defined in the same way; for these the leading part is a smooth section over
$V\setminus0$, homogeneous of degree $m$, i.e.\ a section over the sphere bundle
$SV=(V\setminus0)/\R^+$ of $N_m\otimes W$, where $N_m$ is the line bundle whose
sections are the functions homogeneous of degree $m$ on $V\setminus0$.

\begin{proposition}\label{manifolds:prop:principal-symbol}
There is a unique linear map
\begin{equation*}
\sigma_m:\Psido m(X;E,F)\longrightarrow\Symq m(T^*X;\hom(E,F))
\end{equation*}
such that for every trivializing chart $\kappa:W\to\kappa(W)$ and every
$\psi\in\CIc(W)$, $\sigma_m(\kappa_*(\psi A\psi))$ is the local expression of
$\psi^2\sigma_m(A)$, using the coordinates \eqref{manifolds:eq:cotangent-change}
on $T^*W$. It has the following properties.
\begin{enumerate}[label=\textup{(\alph*)}]
\item The sequence
$0\to\Psido{m-1}(X;E,F)\hookrightarrow\Psido m(X;E,F)\xrightarrow{\sigma_m}\Symq m(T^*X;\hom(E,F))\to0$
is exact.
\item $\sigma_{m+m'}(AB)=\sigma_m(A)\sigma_{m'}(B)$ and $\sigma_m(A^*)=\sigma_m(A)^*$.
\item For $P\in\Diff k(X;E,F)$, $\sigma_k(P)$ is the class of the polynomial
symbol defined by \eqref{manifolds:eq:oscillatory-test}.
\item For classical operators $\sigma_m(A)$ is a homogeneous section and there is a
short exact sequence
\begin{equation*}
\Psido{m-1}_{\mathrm{ph}}(X;E,F)\hookrightarrow\Psido m_{\mathrm{ph}}(X;E,F)
\xrightarrow{\sigma_m}\CI(S^*X;N_m\otimes\hom(E,F))\to0.
\end{equation*}
\end{enumerate}
\end{proposition}
\begin{proof}
Choose a finite cover by trivializing charts $\kappa_i:W_i\to\kappa_i(W_i)$ and
$\psi_i\in\CIc(W_i)$ with $\sum_i\psi_i^2=1$ (e.g.\ $\psi_i=\rho_i(\sum_k\rho_k^2)^{-1/2}$
for a partition of unity $\rho_i$). The principal symbol of
$\kappa_{i*}(\psi_iA\psi_i)\in\Psido m(\R^n)$ has a representative supported, in $x$,
in a compact subset of $\kappa_i(W_i)$, well defined modulo $S^{m-1}$ symbols with
the same property; transferring it to $T^*W_i$ by \eqref{manifolds:eq:cotangent-change}
and extending by zero gives $a_i\in\Symq m(T^*X;\hom(E,F))$. Set
$\sigma_m(A)=\sum_ia_i$.

To check the required property, let $\kappa:W\to\kappa(W)$ and $\psi\in\CIc(W)$ be
arbitrary. Then $\psi A\psi=\sum_i\psi_i^2\psi A\psi$ and
$\psi_i^2(\psi A\psi)-\psi_i(\psi A\psi)\psi_i=\psi_i[\psi_i,\psi A\psi]$, which in
the chart $\kappa$ has principal symbol zero by \ref{manifolds:E3}. The operator
$\psi_i\psi A\psi\psi_i=\psi(\psi_iA\psi_i)\psi$ has kernel supported in
$(W\cap W_i)^2$; in the chart $\kappa_i$ its principal symbol is $\psi^2$ times that
of $\kappa_{i*}(\psi_iA\psi_i)$, and by \ref{manifolds:E7} its principal symbol in
the chart $\kappa$ is the same invariant function on $T^*X$, namely $\psi^2a_i$.
Summing over $i$ gives $\sigma_m(\kappa_*(\psi A\psi))=\psi^2\sigma_m(A)$. Uniqueness
follows by applying the property to the $\psi_i$: two such maps differ by $b$ with
$\psi_i^2b\in S^{m-1}$ for all $i$, so $b=\sum_i\psi_i^2b\in S^{m-1}$.

(a) Clearly $\sigma_m$ vanishes on $\Psido{m-1}$. If $\sigma_m(A)=0$ then for every
trivializing chart and $\psi\in\CIc(W)$, $\kappa_*(\psi A\psi)$ has vanishing
principal symbol, so lies in $\Psido{m-1}(\R^n)$ by \ref{manifolds:E4}; by
\cref{manifolds:thm:local}, $A\in\Psido{m-1}(X;E,F)$. For surjectivity define, with
the notation above, the \emph{quantization map}
\begin{equation}\label{manifolds:eq:quantization}
q(a)=\sum_i\kappa_i^*\circ\big(\tilde\chi_i\Op(a_i)\tilde\chi_i\big)\circ(\kappa_i^{-1})^*,
\end{equation}
where $a_i$ is the local expression of $\psi_i^2a$ in the chart $\kappa_i$ (extended
by zero outside $\kappa_i(W_i)$) and $\tilde\chi_i\in\CIc(\kappa_i(W_i))$ is equal to
one near the support of $\psi_i\circ\kappa_i^{-1}$. By
\cref{manifolds:cor:local-ops}, $q(a)\in\Psido m(X;E,F)$, and applying the defining
property in an arbitrary chart, using \ref{manifolds:E7} as above, gives
$\sigma_m(q(a))=\sum_i\psi_i^2a=a$. Note that $q$ is linear and maps
$S^{-\infty}$ into $\Psido{-\infty}$.

(b) Let $\psi\in\CIc(W)$ and $\chi\in\CIc(W)$ equal to one near $\supp\psi$. Then
$\psi AB\psi-(\psi A\chi)(\chi B\psi)=\psi A(1-\chi^2)B\psi$ is smoothing, by the
argument in the proof of \cref{manifolds:thm:composition}, and the local
representatives of $\psi A\chi$ and $\chi B\psi$ have principal symbols
$\psi\chi\sigma_m(A)$ and $\chi\psi\sigma_{m'}(B)$ (compare $\psi A\chi$ with
$\tilde\psi A\tilde\psi$ for $\tilde\psi=1$ on $\supp\chi$, using \ref{manifolds:E2}).
So \ref{manifolds:E3} gives the multiplicativity. The statement for adjoints was
proved in \cref{manifolds:prop:adjoint}.

(c) is immediate from \eqref{manifolds:eq:diffop}, and (d) follows from (a),
since the quantization of a classical symbol is classical.
\end{proof}

\begin{remark}\label{manifolds:rem:high-road}
Since $\Psido m(X)=\Ic m(X^2,\Diag;\Omega_R)$ is a $\CI(X^2)$-module,
\eqref{manifolds:eq:def-psido} may be written
$\Psido m(X;E,F)=\Psido m(X)\otimes_{\CI(X^2)}\CI(X^2;\Hom(E,F))$. Thus an operator
between sections of bundles is a finite sum $\sum_iA_iG_i$ with $A_i\in\Psido m(X)$
and $G_i\in\CI(X^2;\Hom(E,F))$, and in local trivializations it is simply a matrix
of scalar operators.
\end{remark}

\begin{proposition}\label{manifolds:prop:asymptotic}
If $A_j\in\Psido{m-j}(X;E,F)$ for $j\in\N_0$, there exists $A\in\Psido m(X;E,F)$,
unique modulo $\Psido{-\infty}(X;E,F)$, such that
$A-\sum_{j<N}A_j\in\Psido{m-N}(X;E,F)$ for all $N$; we write $A\sim\sum_jA_j$.
\end{proposition}
\begin{proof}
Choose a finite cover by trivializing charts $U_c$, decompose each $A_j$ as in
\cref{manifolds:lem:localization} with respect to it,
$A_j=\sum_cA_{j,c}+A_j'$, and note that the supports of the kernels of the $A_{j,c}$
may be taken to lie in a fixed compact subset $K_c$ of $U_c\times U_c$ independent of
$j$ (use the same partition of unity for all $j$). By \ref{manifolds:E5} there are
$B_c\in\Psido m(X;E,F)$ with kernels supported in $U_c\times U_c$ and
$B_c-\sum_{j<N}A_{j,c}\in\Psido{m-N}$ for all $N$. Then $A=\sum_cB_c$ satisfies
$A-\sum_{j<N}A_j=\sum_c(B_c-\sum_{j<N}A_{j,c})-\sum_{j<N}A_j'\in\Psido{m-N}$. If $A$ and
$\tilde A$ both have the property, $A-\tilde A\in\bigcap_N\Psido{m-N}=\Psido{-\infty}$.
\end{proof}

Thus $\Psido{}(X;E)=\bigcup_m\Psido m(X;E)$ is a filtered $*$-algebra, the
principal symbol map is a homomorphism onto the symbol algebra, and
$\Psido{-\infty}(X;E)$ is an ideal. Moreover, from \ref{manifolds:C3} or directly
from \cref{manifolds:thm:composition},
\begin{equation*}
\Diff k(X;F,G)\circ\Psido m(X;E,F)\subset\Psido{k+m}(X;E,G),\qquad
\sigma_{k+m}(PA)=\sigma_k(P)\sigma_m(A).
\end{equation*}

\section{Ellipticity and parametrices}\label{manifolds:sec:elliptic}

\begin{definition}\label{manifolds:def:elliptic}
An operator $A\in\Psido m(X;E,F)$ is \emph{elliptic} if there is a symbol
$b\in S^{-m}(T^*X;\hom(F,E))$ such that, for a representative $a$ of $\sigma_m(A)$,
\begin{equation*}
ba-\Id_E\in S^{-1}(T^*X;\hom(E)),\qquad ab-\Id_F\in S^{-1}(T^*X;\hom(F)).
\end{equation*}
\end{definition}

\begin{lemma}\label{manifolds:lem:elliptic}
Fix a metric on $T^*X$. Then $A\in\Psido m(X;E,F)$ is elliptic if and only if a
(any) representative $a$ of $\sigma_m(A)$ is invertible for $|\xi|\ge R$, for some $R$,
with $|a(x,\xi)^{-1}|\le C\ang\xi^{-m}$ there. For classical operators this holds if
and only if the homogeneous principal symbol is invertible on $T^*X\setminus0$, and
for differential operators it is the condition given in
\cref{manifolds:sec:definition}.
\end{lemma}
\begin{proof}
If $ba=\Id+r$, $ab=\Id+r'$ with $r,r'\in S^{-1}$ then for $|\xi|$ large $|r|,|r'|<\frac12$,
so $a$ has a left and a right inverse, hence is invertible (in particular the
ranks of $E$ and $F$ coincide), and $a^{-1}=(\Id+r)^{-1}b$ satisfies the bound.
Conversely, if the bound holds, let $\chi\in\CI(T^*X)$ be equal to one for $|\xi|\le R+1$ and
vanish for $|\xi|\ge R+2$, and set $b=(1-\chi)a^{-1}$. Differentiating
$a\,a^{-1}=\Id$ repeatedly, each derivative of $a^{-1}$ is a sum of products of
$a^{-1}$ with derivatives of $a$, and the symbol estimates for $b\in S^{-m}$
follow by induction; $ab-\Id=-\chi\Id$ and $ba-\Id$ have compact support in $\xi$, so
lie in $S^{-\infty}$. For a homogeneous symbol invertible on the sphere bundle the
bound follows from compactness of $S^*X$ and homogeneity.
\end{proof}

\begin{theorem}\label{manifolds:thm:parametrix}
If $A\in\Psido m(X;E,F)$ is elliptic, there exists $B\in\Psido{-m}(X;F,E)$, unique
modulo $\Psido{-\infty}$, such that
\begin{equation}\label{manifolds:eq:parametrix}
BA=\Id_E-R_E,\quad AB=\Id_F-R_F,\qquad R_E\in\Psido{-\infty}(X;E),\ R_F\in\Psido{-\infty}(X;F).
\end{equation}
If $A$ is classical, so is $B$.
\end{theorem}
\begin{proof}
With $b$ as in \cref{manifolds:def:elliptic} let $B_0=q(b)\in\Psido{-m}(X;F,E)$.
By \cref{manifolds:prop:principal-symbol},
$AB_0=\Id_F-S$ with $S\in\Psido{-1}(X;F)$. By
\cref{manifolds:prop:asymptotic} choose $B_r\in\Psido{-m}(X;F,E)$ with
$B_r\sim\sum_{j\ge0}B_0S^j$. Then for each $N$,
\begin{equation*}
AB_r\equiv AB_0\sum_{j<N}S^j=(\Id-S)\sum_{j<N}S^j=\Id-S^N\mod\Psido{-N},
\end{equation*}
so $AB_r-\Id\in\Psido{-N}$ for all $N$, i.e.\ $AB_r=\Id-R_F$ with $R_F$ smoothing. In
the same way, starting from $B_0A=\Id-S'$, there is $B_l$ with $B_lA=\Id-R_E$. Then
\begin{equation*}
B_l=B_l(AB_r+R_F)=(\Id-R_E)B_r+B_lR_F=B_r-R_EB_r+B_lR_F,
\end{equation*}
so $B_l-B_r\in\Psido{-\infty}$ by \cref{manifolds:cor:smoothing-ideal}, and
$B=B_r$ satisfies both identities in \eqref{manifolds:eq:parametrix}. The same
computation shows that any left parametrix and any right parametrix differ by a
smoothing operator, which gives uniqueness. In the classical case each step stays
in the classical calculus.
\end{proof}

\begin{corollary}[Elliptic regularity]\label{manifolds:cor:regularity}
If $A\in\Psido m(X;E,F)$ is elliptic, $u\in\CmI(X;E)$ and $Au\in\CI(X;F)$, then
$u\in\CI(X;E)$. In particular the null space of $A$ on distributions,
$\Nul(A)=\{u\in\CmI(X;E);Au=0\}$, consists of smooth sections.
\end{corollary}
\begin{proof}
By \eqref{manifolds:eq:parametrix}, $u=BAu+R_Eu$; here $BAu\in\CI$ by
\cref{manifolds:prop:continuity} and $R_Eu\in\CI$ since $R_E$ is smoothing.
\end{proof}

The parametrix reduces the study of $A$ to that of operators $\Id+R$ with $R$
smoothing: $\Nul(A)\subset\Nul(\Id-R_E)$ and $\Ran(A)\supset\Ran(\Id-R_F)$. In
\cref{manifolds:sec:fredholm} we use this to show that elliptic operators are
Fredholm, on smooth sections as well as on Sobolev spaces. The \emph{index}
of an elliptic operator is then
\begin{equation}\label{manifolds:eq:index}
\ind(A)=\dim\Nul(A)-\dim\big(\CI(X;F)/A\,\CI(X;E)\big).
\end{equation}
For a linear map between finite-dimensional spaces the corresponding integer is
just the difference of the dimensions of the spaces; the index is a genuinely
infinite-dimensional invariant. Its computation in terms of the principal symbol
is the content of the Atiyah--Singer index theorem, the subject of
\cref{ch:index} and the following chapters.

\section{Sobolev spaces}\label{manifolds:sec:sobolev}

\subsection{Local facts}
\begin{lemma}\label{manifolds:lem:sobolev-local}
Let $K\subset\R^n$ be compact and $\Hs s_K=\{u\in\Hs s(\R^n);\supp u\subset K\}$.
\begin{enumerate}[label=\textup{(\alph*)}]
\item If $\Phi:\Omega\to\Omega'$ is a diffeomorphism of open sets and $K\subset\Omega$, then
$\Phi_*=(\Phi^{-1})^*:\Hs s_K\to\Hs s_{\Phi(K)}$ is bounded, for every $s\in\R$.
\item An operator with kernel in $\CIc(\R^{2n})$ is bounded $\Hs t(\R^n)\to\Hs{t'}(\R^n)$
for all $t,t'$.
\item \textup{(Rellich)} If $t<s$ the inclusion $\Hs s_K\hookrightarrow\Hs t(\R^n)$ is compact.
\item $\CmIc(\R^n)=\bigcup_s\Hs s_{\mathrm c}(\R^n)$, and $\Hs s(\R^n)\subset\mathcal C^k(\R^n)$,
with bounded inclusion into bounded $\mathcal C^k$ functions, if $s>k+\frac n2$.
\end{enumerate}
\end{lemma}
\begin{proof}
(a) Let $\Lambda_t=\Op(\ang\xi^t)$ and choose $\chi,\chi'\in\CIc(\Omega)$ with $\chi=1$ near
$K$ and $\chi'=1$ near $\supp\chi$. For $u\in\Hs s_K$ put $f=\Lambda_su\in L^2$; then
\begin{equation*}
u=\chi\Lambda_{-s}f=\chi\Lambda_{-s}\chi'f+\chi\Lambda_{-s}(1-\chi')f.
\end{equation*}
The kernel of $\Lambda_{-s}$ is $k(x-y)$ where $k$ is smooth outside $0$ and,
together with all its derivatives, rapidly decreasing at infinity (since
$|z|^{2N}D^\alpha k=\cF^{-1}(\Delta_\xi^N(\xi^\alpha\ang\xi^{-s}))$ is bounded when $2N>n-s+|\alpha|$).
So the second term is given by a smooth kernel $\chi(x)k(x-y)(1-\chi'(y))$, supported
in $\supp\chi\times\R^n$ and with $x$-derivatives bounded by $C_N\ang y^{-N}$; it maps
$L^2$ boundedly into smooth functions supported in $\supp\chi$, with each
$\mathcal C^k$ norm bounded, and these are preserved by $\Phi_*$. For the first term,
$\Phi_*(\chi\Lambda_{-s}\chi'f)=\Phi_*(\chi\Lambda_{-s}\chi')\,\Phi_*(\chi'f)$, where
$\Phi_*(\chi\Lambda_{-s}\chi')\in\Psido{-s}(\R^n)$ by \ref{manifolds:E7}, which is bounded
$L^2\to\Hs s$ by \ref{manifolds:E6}, and $\|\Phi_*(\chi'f)\|_{L^2}\le C\|f\|_{L^2}$ since the
Jacobian of $\Phi^{-1}$ is bounded on $\Phi(\supp\chi')$.

(b) Let $T$ have kernel $k\in\CIc(\R^{2n})$. Then for each multi-index $\alpha$ the
derivative $D^\alpha Tu(x)=u(D^\alpha_xk(x,\cdot))$ is bounded by the product
$\|u\|_{\Hs t}\sup_x\|D_x^\alpha k(x,\cdot)\|_{\Hs{-t}}$ and $Tu$ has support in a
fixed compact set.

(c) Let $u_j\in\Hs s_K$, $\|u_j\|_{\Hs s}\le1$, and $\chi\in\CIc(\R^n)$ equal to one near $K$.
Then $\hat u_j(\xi)=u_j(\chi e^{-ix\cdot\xi})$, so $\hat u_j$ and all its derivatives are
bounded on compact sets uniformly in $j$. By the Arzel\`a--Ascoli theorem and a
diagonal argument a subsequence (still denoted $u_j$) has $\hat u_j$ converging
uniformly on compact sets. Then for any $R$,
\begin{equation*}
\|u_j-u_k\|_{\Hs t}^2\le\int_{|\xi|\le R}\ang\xi^{2t}|\hat u_j-\hat u_k|^2\dbar\xi
+\ang R^{2(t-s)}\int\ang\xi^{2s}|\hat u_j-\hat u_k|^2\dbar\xi ,
\end{equation*}
where the last term is at most $4\ang R^{2(t-s)}$; choosing $R$ large and then $j,k$
large shows that $(u_j)$ is Cauchy in $\Hs t$.

(d) A compactly supported distribution has finite order, so $|\hat u(\xi)|\le C\ang\xi^N$
for some $N$ and $u\in\Hs{-N-n}$. If $s>k+\frac n2$ and $|\alpha|\le k$ then
$\int|\xi^\alpha\hat u|\,\dbar\xi\le\|\ang\xi^{|\alpha|-s}\|_{L^2}\|u\|_{\Hs s}$, so
$D^\alpha u$ is continuous and bounded.
\end{proof}

\subsection{Sobolev spaces on a compact manifold}
Fix a finite cover of $X$ by charts $\kappa_i:W_i\to\kappa_i(W_i)$ over which $E$ is
trivial, and a subordinate partition of unity $\rho_i$.

\begin{definition}\label{manifolds:def:sobolev}
For $s\in\R$,
\begin{equation*}
\Hs s(X;E)=\{u\in\CmI(X;E);\ \kappa_{i*}(\rho_iu)\in\Hs s(\R^n)^{N_E}\ \forall i\},
\qquad\|u\|_s^2=\sum_i\|\kappa_{i*}(\rho_iu)\|_{\Hs s}^2 .
\end{equation*}
\end{definition}

\begin{proposition}\label{manifolds:prop:sobolev}
\begin{enumerate}[label=\textup{(\alph*)}]
\item $\Hs s(X;E)$ is independent of the choices made, with equivalent norms. A
distribution $u$ lies in $\Hs s(X;E)$ if and only if $\kappa_*(\psi u)\in\Hs s(\R^n)^{N_E}$
for every trivializing chart $\kappa$ and $\psi\in\CIc(W)$; and if $w$ has support in a
compact subset of $W$ and $\kappa_*w\in\Hs s$, then $w\in\Hs s(X;E)$, with
$\|w\|_s\le C\|\kappa_*w\|_{\Hs s}$, $C$ depending only on the support.
\item $\Hs s(X;E)$ is a Hilbert space in which $\CI(X;E)$ is dense, and
$\Hs 0(X;E)=L^2(X;E)$ with equivalent norms.
\item $\bigcap_s\Hs s(X;E)=\CI(X;E)$ and $\bigcup_s\Hs s(X;E)=\CmI(X;E)$.
\item For $t<s$ the inclusion $\Hs s(X;E)\hookrightarrow\Hs t(X;E)$ is compact.
\item The $L^2$ pairing \eqref{manifolds:eq:L2} extends by continuity to a
sesquilinear pairing $\Hs s(X;E)\times\Hs{-s}(X;E)\to\C$, and every continuous linear
functional on $\Hs s(X;E)$ is of the form $u\mapsto\inner uw$ for a unique
$w\in\Hs{-s}(X;E)$.
\end{enumerate}
\end{proposition}
\begin{proof}
(a) Let $\kappa$ be a trivializing chart and $\psi\in\CIc(W)$. Then
$\psi u=\sum_i\psi\rho_iu$, and $\kappa_{i*}(\psi\rho_iu)=(\kappa_{i*}\psi)\kappa_{i*}(\rho_iu)\in\Hs s$
(by \ref{manifolds:E6}, after changing trivializations, which is multiplication by
a smooth matrix) is supported in a compact subset of $\kappa_i(W\cap W_i)$, so by
\cref{manifolds:lem:sobolev-local}(a) applied to $\kappa\circ\kappa_i^{-1}$,
$\kappa_*(\psi\rho_iu)\in\Hs s$, with norm estimates. This proves that membership
for one choice of cover and partition implies the condition for all charts, which in
turn implies membership for any other choice; the norm equivalences come with the
estimates. The last statement follows in the same way.

(b) The map $u\mapsto(\kappa_{i*}(\rho_iu))_i$ is an isometry onto a subspace of
$\bigoplus_i\Hs s(\R^n)^{N_E}$. Its range is closed: if $v_i^{(k)}=\kappa_{i*}(\rho_iu^{(k)})\to v_i$
in $\Hs s$, then $u^{(k)}=\sum_i\kappa_i^*v_i^{(k)}$ converges in $\CmI(X;E)$ to
$u=\sum_i\kappa_i^*v_i$ (the $v_i$ being supported in $\kappa_i(\supp\rho_i)$), and so
$v_i=\kappa_{i*}(\rho_iu)$. For density, approximate each $\kappa_{i*}(\rho_iu)$ in $\Hs s$
by smooth functions supported in a fixed compact subset of $\kappa_i(W_i)$ (by
convolution with an approximate identity) and use (a). The identification of $\Hs0$
with $L^2$ follows since the density $\nu$ and the metric on $E$ are given by smooth,
positive, locally bounded factors.

(c) If $u\in\Hs s$ for all $s$, each $\kappa_{i*}(\rho_iu)$ is smooth by
\cref{manifolds:lem:sobolev-local}(d), so $u=\sum_i\rho_iu$ is smooth. Conversely,
each $\kappa_{i*}(\rho_iu)$ is a compactly supported distribution, so lies in some
$\Hs{s_i}$.

(d) Under the isometry of (b), $\Hs s(X;E)$ maps into
$\bigoplus_i(\Hs s_{K_i})^{N_E}$ with $K_i=\kappa_i(\supp\rho_i)$, and the inclusion of this space
into $\bigoplus_i\Hs t(\R^n)^{N_E}$ is compact by \cref{manifolds:lem:sobolev-local}(c).

(e) Choose $\psi_i\in\CIc(W_i)$ equal to one on $\supp\rho_i$. In local coordinates
$\nu=\nu_i|dx|$ and the metric is $h_i$, so
$\inner uw=\sum_i\int\langle h_i\nu_i\kappa_{i*}(\rho_iu),\kappa_{i*}(\psi_iw)\rangle dx$,
which is bounded by $C\|u\|_s\|w\|_{-s}$ using $|\int f\bar g|\le\|f\|_{\Hs s}\|g\|_{\Hs{-s}}$ on
$\R^n$ and (a). Conversely, let $\ell$ be a continuous linear functional on
$\Hs s(X;E)$. By the Riesz representation theorem, $\ell(u)$ is the $\Hs s$ inner
product of $u$ with some $w'\in\Hs s$, i.e.
\begin{equation*}
\ell(u)=\sum_i\int\big\langle\kappa_{i*}(\rho_iu),\Lambda_{2s}v_i\big\rangle dx
=\sum_i\int\big\langle\kappa_{i*}(\rho_iu),\chi_i\Lambda_{2s}v_i\big\rangle dx,
\end{equation*}
where $v_i=\kappa_{i*}(\rho_iw')$, $\Lambda_{2s}=\Op(\ang\xi^{2s})$ and $\chi_i\in\CIc(\kappa_i(W_i))$
equals one on $\kappa_i(\supp\rho_i)$. Here $\chi_i\Lambda_{2s}v_i\in\Hs{-s}$ has compact support in
$\kappa_i(W_i)$, so $w=\sum_i\rho_i\kappa_i^*\big(\nu_i^{-1}h_i^{-1}\chi_i\Lambda_{2s}v_i\big)\in\Hs{-s}(X;E)$ by (a), and
$\ell(u)=\inner uw$. Uniqueness follows from density of $\CI(X;E)$ in $\Hs s(X;E)$.
\end{proof}

\begin{theorem}\label{manifolds:thm:sobolev-bounded}
Each $A\in\Psido m(X;E,F)$ extends to a bounded operator
$A:\Hs s(X;E)\to\Hs{s-m}(X;F)$ for every $s\in\R$. In particular smoothing operators
map $\Hs s$ into $\Hs t$ for all $s,t$, and are compact on each $\Hs s$; operators of
negative order are compact on each $\Hs s$.
\end{theorem}
\begin{proof}
Use the decomposition \eqref{manifolds:eq:localization} with respect to a
cover by charts trivializing $E$ and $F$. The smoothing part $A'$ maps
$\Hs s(X;E)$ to $\Hs t(X;F)$ for all $t$: decomposing $A'=\sum_{j,k}\rho_jA'\rho_k$, each
term is, in the charts $\kappa_j$ and $\kappa_k$, an operator with kernel in
$\CIc(\R^{2n})$, so \cref{manifolds:lem:sobolev-local}(b) and
\cref{manifolds:prop:sobolev}(a) apply. For $A_i$, with kernel supported in
$K'\times K'$, $K'\Subset W_i$, choose $\psi\in\CIc(W_i)$ equal to one near $K'$. Then
$A_iu=A_i(\psi u)$, $\kappa_{i*}(\psi u)\in\Hs s$, and the local representative of
$A_i$ is a matrix of elements of $\Psido m(\R^n)$, so maps it into $\Hs{s-m}$,
with support in $\kappa_i(K')$; by \cref{manifolds:prop:sobolev}(a),
$A_iu\in\Hs{s-m}(X;F)$ with the appropriate bound. The compactness statements
follow from \cref{manifolds:prop:sobolev}(d), since an operator of order $m<0$ maps
$\Hs s$ boundedly into $\Hs{s-m}$, which includes compactly into $\Hs s$.
\end{proof}

\begin{proposition}\label{manifolds:prop:elliptic-estimate}
Let $A\in\Psido m(X;E,F)$ be elliptic. If $u\in\CmI(X;E)$ and $Au\in\Hs{s-m}(X;F)$ then
$u\in\Hs s(X;E)$, and for any $t<s$
\begin{equation*}
\|u\|_s\le C\big(\|Au\|_{s-m}+\|u\|_t\big).
\end{equation*}
Consequently, for any elliptic $\Lambda\in\Psido s(X;E)$,
\begin{equation*}
\Hs s(X;E)=\{u\in\CmI(X;E);\ \Lambda u\in L^2\}=\{u\in\CmI(X;E);\ Pu\in L^2\ \forall P\in\Psido s(X;E)\}.
\end{equation*}
\end{proposition}
\begin{proof}
By \eqref{manifolds:eq:parametrix}, $u=BAu+R_Eu$, with $BAu\in\Hs s$ by
\cref{manifolds:thm:sobolev-bounded} and $R_Eu\in\CI$ (note $u\in\Hs t$ for some
$t$ by \cref{manifolds:prop:sobolev}(c)); the estimate follows from the bounds on
$B$ and $R_E$. For the last statement, an elliptic $\Lambda\in\Psido s(X;E)$ exists
(e.g.\ $q(\ang\xi^s\Id_E)$ for a metric on $T^*X$), and the inclusions follow from
\cref{manifolds:thm:sobolev-bounded} and the first part.
\end{proof}

\section{Fredholm theory of elliptic operators}\label{manifolds:sec:fredholm}

\begin{theorem}\label{manifolds:thm:fredholm}
Let $A\in\Psido m(X;E,F)$ be elliptic, and let $A^*\in\Psido m(X;F,E)$ be its adjoint.
\begin{enumerate}[label=\textup{(\alph*)}]
\item $N=\Nul(A)$ and $N^*=\Nul(A^*)$ (on distributions) are finite-dimensional
subspaces of $\CI(X;E)$ and $\CI(X;F)$ respectively.
\item For every $s\in\R$, $A:\Hs s(X;E)\to\Hs{s-m}(X;F)$ is Fredholm, with null space
$N$ and range $\{f\in\Hs{s-m}(X;F);\ \inner fv=0\ \forall v\in N^*\}$.
\item $A:\CI(X;E)\to\CI(X;F)$ has null space $N$ and closed range, with complement
$N^*$: $\CI(X;F)=A\,\CI(X;E)\oplus N^*$, an $L^2$-orthogonal sum.
\item $\ind(A)=\dim N-\dim N^*$ is the index on $\CI$ and on every $\Hs s$. It
depends only on $\sigma_m(A)$; moreover $\ind(A^*)=-\ind(A)$, $\ind(AB)=\ind(A)+\ind(B)$
for elliptic $B\in\Psido{m'}(X;G,E)$, and the index is constant on any family of
elliptic operators which is continuous in the norm of
$\Lin(\Hs s(X;E),\Hs{s-m}(X;F))$ for some $s$.
\end{enumerate}
\end{theorem}
\begin{proof}
Let $B$ be a parametrix as in \eqref{manifolds:eq:parametrix}. By
\cref{manifolds:thm:sobolev-bounded}, $A:\Hs s\to\Hs{s-m}$ and $B:\Hs{s-m}\to\Hs s$ are
bounded and $R_E$, $R_F$ are compact on $\Hs s$ and $\Hs{s-m}$, so $A$ is Fredholm on these
spaces by the parametrix characterization of Fredholm operators
(\cref{hilbert:thm:atkinson}). Its null space on $\Hs s$ is contained in $N$, and
$N\subset\CI\subset\Hs s$ by \cref{manifolds:cor:regularity}; so it equals $N$, which is
therefore finite-dimensional. The adjoint is elliptic, since
$\sigma_m(A^*)=\sigma_m(A)^*$ is inverted by $b^*$, so the same applies to $N^*$. This
proves (a).

For (b), let $R_s=A\Hs s(X;E)$, which is closed, and
$Y_s=\{f\in\Hs{s-m};\inner fv=0\ \forall v\in N^*\}$. If $u\in\Hs s$ and $v\in N^*$, then
$\inner{Au}v=\inner u{A^*v}=0$, by continuity from smooth $u$; so $R_s\subset Y_s$.
If $R_s\ne Y_s$ there is a continuous linear functional on $\Hs{s-m}$ which vanishes on
$R_s$ but not on $Y_s$. By \cref{manifolds:prop:sobolev}(e) it is $f\mapsto\inner fw$ for some
$w\in\Hs{m-s}(X;F)$; then $\inner u{A^*w}=\inner{Au}w=0$ for all $u\in\CI(X;E)$, so $A^*w=0$,
$w\in N^*$, and the functional vanishes on $Y_s$, a contradiction.

(c) If $f\in\CI(X;F)$ is orthogonal to $N^*$ then by (b) $f=Au$ with $u\in\Hs s$ for
any $s$, and $u$ is smooth by \cref{manifolds:cor:regularity}. The range on $\CI$ is
thus the orthocomplement of $N^*$ in $\CI(X;F)$, which is closed, and since $N^*\subset\CI$,
$f=(f-\pi_{N^*}f)+\pi_{N^*}f$ gives the decomposition, where $\pi_{N^*}$ is orthogonal
projection onto $N^*$.

(d) By (b), the codimension of the range on $\Hs{s-m}$ is $\dim N^*$, the elements of
$N^*$ being linearly independent functionals on $\Hs{s-m}$; similarly on $\CI$ by
(c). If $B'\in\Psido{m-1}(X;E,F)$, then $B':\Hs s\to\Hs{s-m}$ is compact (it maps into
$\Hs{s-m+1}$), so $\ind(A+B')=\ind(A)$ by the stability of the index under compact
perturbations (\cref{hilbert:thm:stability}). The formula for $A^*$ is clear from
(a), the product formula follows from \cref{hilbert:prop:index-product} applied to
$\Hs{s}\to\Hs{s-m'}\to\Hs{s-m'-m}$, and the constancy along continuous families is
\cref{hilbert:thm:stability} again.
\end{proof}

In particular, for $R\in\Psido{-\infty}(X;E)$ the operator $\Id+R$ is Fredholm of
index zero on $L^2(X;E)$ and on $\CI(X;E)$: it is elliptic of order $0$, and
$\Id+tR$, $t\in[0,1]$, is a continuous family connecting it to the identity. Note
also that the range in (c) is closed in $\CI(X;F)$ and has a finite-dimensional
complement; closedness of the range is part of the Fredholm property, and is not
implied by the finite codimension of its closure.

\subsection{Generalized inverses}
We need a lemma on composites with smoothing operators.

\begin{lemma}\label{manifolds:lem:sandwich}
Let $R_1\in\Psido{-\infty}(X;G',G)$, $R_2\in\Psido{-\infty}(X;E,E')$, and let
$T:\Hs t(X;E')\to\Hs{t'}(X;G')$ be bounded, for some $t,t'\in\R$. Then
$R_1TR_2\in\Psido{-\infty}(X;E,G)$, and its kernel depends continuously, in
$\CI(X^2;\Hom(E,G)\otimes\Omega_R)$, on $T$ in the operator norm.
\end{lemma}
\begin{proof}
The kernel of $R_1TR_2$ is $K(x,z)=\big(R_1\,T\,k_z\big)(x)$, where $k_z=K_{R_2}(\cdot,z)$
(with the fibre of $E'_z\otimes\Omega_z$ trivialized locally in $z$). The map $z\mapsto k_z$
is smooth into $\Hs t(X;E')$, so $z\mapsto Tk_z$ is smooth into $\Hs{t'}(X;G')$ with
derivatives bounded by $\|T\|$ times those of $k_z$, and
$R_1:\Hs{t'}(X;G')\to\CI(X;G)$ is continuous, with $D^\alpha_x(R_1g)(x)$ the pairing of $g$
with the smooth family $D^\alpha_xK_{R_1}(x,\cdot)\in\Hs{-t'}$. Hence $K$ is smooth, and each
$\mathcal C^k$ norm of $K$ is bounded by $C_k\|T\|$.
\end{proof}

\begin{proposition}\label{manifolds:prop:generalized-inverse}
Let $A\in\Psido m(X;E,F)$ be elliptic and let $\pi_N$, $\pi_{N^*}$ be the orthogonal
projections onto $N=\Nul(A)$ and $N^*=\Nul(A^*)$. These are smoothing operators of
finite rank, and there is a unique $G\in\Psido{-m}(X;F,E)$ such that
\begin{equation}\label{manifolds:eq:generalized-inverse}
GA=\Id_E-\pi_N,\qquad AG=\Id_F-\pi_{N^*},\qquad G\pi_{N^*}=0 .
\end{equation}
Moreover $G-B\in\Psido{-\infty}$ for any parametrix $B$.
\end{proposition}
\begin{proof}
If $e_1,\dots,e_k$ is an orthonormal basis of $N\subset\CI(X;E)$, then
$\pi_Nu=\sum_j\inner u{e_j}e_j$ has kernel $\sum_je_j(x)\otimes\langle\cdot,e_j(y)\rangle\nu(y)$,
which is smooth; similarly for $\pi_{N^*}$. For $f\in\Hs{s-m}(X;F)$ let $Gf$ be the unique
$u\in\Hs s(X;E)$ with $u\perp N$ and $Au=(\Id-\pi_{N^*})f$; it exists by
\cref{manifolds:thm:fredholm}(b), and does not depend on $s$. The map $A$ restricted to
$N^\perp\cap\Hs s$ is a continuous bijection onto the closed range, so $G$ is bounded
$\Hs{s-m}\to\Hs s$ by the open mapping theorem. The identities
\eqref{manifolds:eq:generalized-inverse} follow directly from the definition, and they
determine $G$. Now
\begin{align*}
G&=G(AB+R_F)=(\Id-\pi_N)B+GR_F,\\
GR_F&=(BA+R_E)GR_F=B(\Id-\pi_{N^*})R_F+R_EGR_F .
\end{align*}
Here $R_EGR_F$ is smoothing by \cref{manifolds:lem:sandwich}, and the remaining terms
other than $B$ are smoothing by \cref{manifolds:cor:smoothing-ideal}. Thus
$G-B\in\Psido{-\infty}$.
\end{proof}

The generalized inverse is the unique linear map from $\CI(X;F)$ to $\CI(X;E)$ which
vanishes on $N^*$, has range $N^\perp$ (the orthocomplement in $\CI(X;E)$) and inverts
$A$ as a map from $N^\perp$ to the range of $A$.

\begin{corollary}\label{manifolds:cor:invertible}
For each $s\in\R$ and each bundle $E$ there exists an invertible, formally
self-adjoint, elliptic $\Lambda_s\in\Psido s(X;E)$ with
$\Lambda_s^{-1}\in\Psido{-s}(X;E)$ and $\inner{\Lambda_su}u\ge0$. In particular
$\|\Lambda_su\|_{L^2}$ is a norm on $\Hs s(X;E)$ equivalent to $\|u\|_s$, and
$\Lambda_s:\Hs t\to\Hs{t-s}$ is an isomorphism for every $t$.
\end{corollary}
\begin{proof}
Let $Q\in\Psido{s/2}(X;E)$ be elliptic, and $\Lambda_s=Q^*Q+\pi$, where $\pi$ is the
orthogonal projection onto the finite-dimensional space $\Nul(Q)\subset\CI$. Then $\Lambda_s$ is
self-adjoint and elliptic, with principal symbol $\sigma(Q)^*\sigma(Q)$. If
$\Lambda_su=0$ then $u$ is smooth and
$0=\inner{\Lambda_su}u=\|Qu\|^2+\|\pi u\|^2$, so $u\in\Nul(Q)$ and $u=\pi u=0$. Thus
$\Nul(\Lambda_s)=\Nul(\Lambda_s^*)=0$, and $\Lambda_s$ is invertible on $\CI$ and on every
$\Hs t$ with inverse its generalized inverse, an element of $\Psido{-s}$.
\end{proof}

\begin{proposition}\label{manifolds:prop:perturbation}
Let $A\in\Psido m(X;E,F)$ be elliptic. There exists a finite-rank
$R\in\Psido{-\infty}(X;E,F)$ such that $A+\eps R$ is invertible for all $\eps\ne0$ if and only if
$\ind(A)=0$. If $A+R$ is invertible, its inverse lies in $\Psido{-m}(X;F,E)$.
\end{proposition}
\begin{proof}
If $A+\eps R$ is invertible for some $\eps$ then $\ind(A)=\ind(A+\eps R)=0$, by
\cref{manifolds:thm:fredholm}(d). Conversely, if $\ind(A)=0$, choose orthonormal bases
$e_j$ of $N$ and $f_j$ of $N^*$ (of the same dimension) and set
$Ru=\sum_j\inner u{e_j}f_j$. If $(A+\eps R)u=0$ then $Au=-\eps Ru\in\Ran(A)\cap N^*=\{0\}$, so
$u\in N$ and $Ru=0$, hence $u=0$. Thus $A+\eps R$ is injective, of index zero, hence
surjective, and its inverse is its generalized inverse, an element of
$\Psido{-m}$.
\end{proof}

\begin{theorem}\label{manifolds:thm:eigen}
Let $A\in\Psido m(X;E)$, $m>0$, be elliptic and formally self-adjoint. Then
$L^2(X;E)$ has an orthonormal basis of eigenfunctions $Ae_j=\lambda_je_j$ with
$e_j\in\CI(X;E)$, $\lambda_j\in\R$, each eigenvalue of finite multiplicity, and
$|\lambda_j|\to\infty$.
\end{theorem}
\begin{proof}
The generalized inverse $G\in\Psido{-m}(X;E)$ of $A$ is self-adjoint (the
characterizing identities \eqref{manifolds:eq:generalized-inverse} are preserved by
taking adjoints, since $N=N^*$) and compact on $L^2$ by
\cref{manifolds:thm:sobolev-bounded}. By the spectral theorem for compact
self-adjoint operators (\cref{hilbert:thm:spectral}), $L^2$ is the orthogonal sum of
$\Nul(G)=N$ and the eigenspaces of $G$ for non-zero eigenvalues $\mu_j\to0$, each
finite-dimensional. If $Gu=\mu u$ with $\mu\ne0$ then $u\perp N$ and
$u=AGu/\mu=Au/\mu$, so $Au=\mu^{-1}u$ and $u$ is smooth by elliptic regularity. Taking
orthonormal bases of $N$ (eigenvalue $0$) and of these eigenspaces gives the result,
with $\lambda_j=\mu_j^{-1}$.
\end{proof}

\section{Smoothing operators}\label{manifolds:sec:smoothing}

Most results about elliptic operators are ultimately reduced, via the parametrix,
to properties of smoothing operators. We fix a smooth positive density $\nu$ and
hermitian metrics, and write the kernel of $A\in\Psido{-\infty}(X;E,F)$ as
$A(x,y)=k_A(x,y)\nu(y)$ with $k_A\in\CI(X^2;\Hom(E,F))$. By the Cauchy--Schwarz
inequality,
\begin{equation}\label{manifolds:eq:HS-bound}
\|Au\|_{L^2}\le\|k_A\|_{L^2(X^2)}\|u\|_{L^2},
\end{equation}
and indeed $A$ is Hilbert--Schmidt with $\|A\|_{2}=\|k_A\|_{L^2(X^2)}$
(\cref{hilbert:prop:HS-kernel}); in particular it is compact on $L^2$.

\begin{lemma}\label{manifolds:lem:HS}
If $A\in\Psido m(X;E,F)$ with $m<-n/2$ then $A$ is Hilbert--Schmidt on $L^2$.
\end{lemma}
\begin{proof}
By \cref{manifolds:lem:localization} it suffices to consider an operator given in
local coordinates by $\Op(a)$, $a\in S^m$, with kernel of compact support. Its kernel
is $K(x,x-z)=\cF^{-1}_{\xi\to z}a(x,\cdot)(z)$, so by Plancherel
$\int\!\!\int|K(x,y)|^2dy\,dx=\int\!\!\int|a(x,\xi)|^2\dbar\xi\,dx$ over the compact $x$-support, which is
finite since $\ang\xi^{2m}$ is integrable.
\end{proof}

\subsection{The group $\Gsm$}
\begin{lemma}\label{manifolds:lem:inverse-smoothing}
If $A\in\Psido{-\infty}(X;E)$ and $\Id+A$ is invertible on $L^2(X;E)$, then
$(\Id+A)^{-1}=\Id+B$ with $B\in\Psido{-\infty}(X;E)$. In particular this holds if
$\|A\|_{\Lin(L^2)}<1$, when $B=\sum_{l\ge1}(-A)^l$.
\end{lemma}
\begin{proof}
Let $T=(\Id+A)^{-1}$. From $T=\Id-AT=\Id-TA$ we get $T=\Id-A(\Id-TA)=\Id-A+ATA$, so
$B=-A+ATA$, which is smoothing by \cref{manifolds:lem:sandwich}.
\end{proof}

\begin{definition}\label{manifolds:def:G}
For a compact manifold $X$ and a vector bundle $E$ over it,
\begin{equation*}
\begin{aligned}
\Gsm(X;E)&=\{A\in\Psido{-\infty}(X;E);\ \Id+A \text{ is invertible on } L^2(X;E)\}\\
&=\{A;\ (\Id+A)^{-1}=\Id+B,\ B\in\Psido{-\infty}(X;E)\},
\end{aligned}
\end{equation*}
topologized as a subset of the Fr\'echet space
$\Psido{-\infty}(X;E)=\CI(X^2;\Hom(E)\otimes\Omega_R)$, with the group law
$(A,B)\mapsto A+B+AB$, i.e.\ composition of $\Id+A$ and $\Id+B$.
\end{definition}

By \cref{manifolds:thm:fredholm}, $\Id+A$ is invertible on $L^2$ if and only if it is
invertible on $\CI(X;E)$, i.e.\ is injective there, since the index is zero.

\begin{proposition}\label{manifolds:prop:G}
$\Gsm(X;E)$ is an open subset of $\Psido{-\infty}(X;E)$ and a topological group.
\end{proposition}
\begin{proof}
Let $A\in\Gsm$ and $\delta=\|(\Id+A)^{-1}\|^{-1}$. If $\|B\|_{\Lin(L^2)}<\delta$ then
$\Id+A+B=(\Id+A)(\Id+(\Id+A)^{-1}B)$ is invertible; by \eqref{manifolds:eq:HS-bound}
this holds for $B$ in a neighbourhood of zero in $\Psido{-\infty}$, so $\Gsm$ is open.
Composition, with kernel $k_A+k_B+\int k_A(x,y)k_B(y,z)\nu(y)$, is continuous; the inverse
$B=-A+A(\Id+A)^{-1}A$ depends continuously on $A$ by \cref{manifolds:lem:sandwich},
since $A\mapsto(\Id+A)^{-1}$ is continuous into $\Lin(L^2)$.
\end{proof}

One should think of $\Gsm(X;E)$ as the group of invertible operators on $L^2(X;E)$
which differ from the identity by an operator with smooth kernel; it plays the
role of an infinite-dimensional general linear group, as the following results
show.

\subsection{Finite-rank approximation}
\begin{lemma}\label{manifolds:lem:finite-rank}
A smoothing operator $A\in\Psido{-\infty}(X;E,F)$ has finite rank if and only if
\begin{equation}\label{manifolds:eq:finite-rank}
A(x,y)=\sum_{i=1}^Nf_i(x)e_i(y)\nu(y),\qquad f_i\in\CI(X;F),\ e_i\in\CI(X;E').
\end{equation}
\end{lemma}
\begin{proof}
Such an operator has range in the span of the $f_i$. Conversely, if the range is
finite-dimensional, it lies in $\CI(X;F)$; let $f_1,\dots,f_N$ be an orthonormal basis of
it. Then $Au=\sum_i\inner{Au}{f_i}f_i$ and
$\inner{Au}{f_i}=\inner u{A^*f_i}=\int_X(A^*f_i)^\dagger(y)u(y)\nu(y)$, where
$(A^*f_i)^\dagger\in\CI(X;E')$ is the dual of $A^*f_i$ under the metric; this is
\eqref{manifolds:eq:finite-rank} with $e_i=(A^*f_i)^\dagger$.
\end{proof}

If the $e_i$ and $f_j$ in \eqref{manifolds:eq:finite-rank} are taken to be orthonormal,
and $E=F$, one may orthonormalize all of them together, obtaining a finite
orthonormal set $v_1,\dots,v_M\subset\CI(X;E)$ and
$Au=\sum_{j,k}a_{jk}\inner u{v_k}v_j$; the finite-rank smoothing operators with range
and coefficients in a fixed such span form a subalgebra isomorphic, as an algebra,
to $M(M,\C)$.

\begin{lemma}\label{manifolds:lem:density}
The finite-rank operators are dense in $\Psido{-\infty}(X;E,F)$. More precisely,
for every $A$ there are finite-rank $A_N$ with $k_{A_N}\to k_A$ in $\CI(X^2;\Hom(E,F))$,
and for a smooth family $A(p)$ depending on parameters $p$ in a compact manifold $P$
the $A_N(p)$ may be chosen as in \eqref{manifolds:eq:finite-rank} with $f_i$, $e_i$
independent of $p$ and coefficients depending smoothly on $p$, converging
uniformly with all derivatives.
\end{lemma}
\begin{proof}
Consider first the torus $\mathbb T^n=\R^n/2\pi\Z^n$ with density $|d\theta|$ and trivial
bundles. A kernel $a\in\CI(\mathbb T^{2n})$ has a rapidly convergent Fourier expansion
\begin{equation*}
a(\theta,\theta')=\sum_{I,J\in\Z^n}a_{IJ}e^{iI\cdot\theta}e^{-iJ\cdot\theta'},\quad
a_{IJ}=(2\pi)^{-2n}\int_{\mathbb T^{2n}}e^{-iI\cdot\theta+iJ\cdot\theta'}a(\theta,\theta')\,d\theta\,d\theta',
\end{equation*}
the coefficients being rapidly decreasing in $(I,J)$ since $a$ is smooth. Hence the
truncations $a_N=\sum_{|I|,|J|\le N}a_{IJ}e^{iI\cdot\theta}e^{-iJ\cdot\theta'}$, which are
finite-rank kernels, converge to $a$ in $\CI$, and continuously in parameters.

In general choose a finite cover of $X$ by charts $\kappa_p:U_p\to\kappa_p(U_p)$ with
bounded images, over which $E$ and $F$ are trivial, and a subordinate partition of unity
$\rho_p$. Then $A=\sum_{p,q}\rho_pA\rho_q$. In the coordinates $\kappa_p(x)$, $\kappa_q(y)$
(rescaled and translated so that $\kappa_p(U_p)$ and $\kappa_q(U_q)$ lie in $(0,2\pi)^n$) the
kernel of $\rho_pA\rho_q$, with respect to the coordinate density, is a matrix of
functions $a\in\CIc((0,2\pi)^{2n})$, which we regard as functions on $\mathbb T^{2n}$. Choose
$\chi_p\in\CIc(\kappa_p(U_p))$ equal to one near $\kappa_p(\supp\rho_p)$, similarly $\chi_q$.
Then $\chi_p(\theta)a_N(\theta,\theta')\chi_q(\theta')$ converges to $a$ in $\CI$, and is a finite
sum of products of functions of $\theta$ and of $\theta'$, compactly supported in the
coordinate patches. Transferring back to $X$ and summing over $p,q$ gives finite-rank
smoothing operators converging to $A$; the functions involved depend only on $N$,
not on parameters.
\end{proof}

\begin{proposition}\label{manifolds:prop:G-connected}
$\Gsm(X;E)$ is connected (by smooth paths).
\end{proposition}
\begin{proof}
Let $A\in\Gsm$ and $\delta=\|(\Id+A)^{-1}\|^{-1}_{\Lin(L^2)}$. As in the proof of
\cref{manifolds:prop:G}, the convex set $\{A+B;\ \|B\|_{\Lin(L^2)}<\delta\}\cap\Psido{-\infty}$
lies in $\Gsm$. By \cref{manifolds:lem:density} and \eqref{manifolds:eq:HS-bound} there
is a finite-rank $A_N$ in this set, and the segment from $A$ to $A_N$ lies in $\Gsm$.
Let $V\subset\CI(X;E)$ be a finite-dimensional space containing the range of $A_N$ and the
range of $A_N^*$, with orthonormal basis $v_1,\dots,v_M$. Then $A_N$ vanishes on $V^\perp$ and
maps $V$ to itself, with matrix $a$, so $\Id+A_N=(\Id_V+a)\oplus\Id_{V^\perp}$ and
$\Id+a\in\GL(M,\C)$. Since $\GL(M,\C)$ is connected, there is a smooth path $g(t)$ in it
from $\Id+a$ to $\Id$, and $\sum_{j,k}(g(t)-\Id)_{jk}\inner\cdot{v_k}v_j$ is a smooth path in
$\Gsm$ from $A_N$ to $0$.
\end{proof}

\subsection{Odd K-theory}
The homotopy classes of maps into $\Gsm$ give a direct definition of odd K-theory:
for a compact manifold $X$ set
\begin{equation}\label{manifolds:eq:K1}
K^{-1}(X)=[X;\Gsm(Y;E)],
\end{equation}
the set of smooth homotopy classes of smooth maps $X\to\Gsm(Y;E)$, where $Y$ is any
compact manifold of positive dimension and $E$ any vector bundle over $Y$. Here a
smooth map is an element
\begin{equation*}
K\in\CI(X\times Y^2;\Hom(E)\otimes\Omega_R)\ \text{ with }\ \Id+K(x,\cdot)\ \text{invertible for every } x\in X,
\end{equation*}
and two such are homotopic if they are the ends of a smooth map
$X\times[0,1]\to\Gsm(Y;E)$. Implicit in \eqref{manifolds:eq:K1} is the following fact;
the group structure is induced by that of $\Gsm$, and the relation with the usual
definition by stable homotopy classes of maps into $\GL(N,\C)$ is visible in the proof.

\begin{proposition}\label{manifolds:prop:K1}
For any two choices $(Y,E)$, $(Z,F)$ with $\dim Y,\dim Z>0$ there is a natural
bijection $[X;\Gsm(Y;E)]\cong[X;\Gsm(Z;F)]$.
\end{proposition}
\begin{proof}[Sketch of proof]
It suffices to compare with the model $Y=\Sphere$, $E=\C$. Embed $\GL(N,\C)$ into
$\Gsm(\Sphere)$ by
\begin{equation*}
(a_{kl})\longmapsto\sum_{k,l=1}^N(a_{kl}-\delta_{kl})\,e^{ik\theta}e^{-il\theta'}\frac{|d\theta'|}{2\pi};
\end{equation*}
the identity is subtracted so that this is a group homomorphism. Given a smooth map
$A:X\to\Gsm(Y;E)$, \cref{manifolds:lem:density} gives finite-rank approximations $A_{(N)}(x)$,
with fixed functions $f_i,e_i$, converging uniformly with all derivatives; for $N$ large
the straight line from $A$ to $A_{(N)}$ stays in $\Gsm$, by the openness estimate,
which is uniform since $X$ is compact. Choosing an orthonormal basis of the span of
the $f_i$ and $e_i$, as in \cref{manifolds:prop:G-connected}, identifies $\Id+A_{(N)}(x)$ with a smooth
map into some $\GL(M,\C)$, and hence, via the embedding, a map $\tilde A:X\to\Gsm(\Sphere)$.
A different basis conjugates $\tilde A$ by a fixed element of $\Gsm(\Sphere)$, which can
be connected to the identity by \cref{manifolds:prop:G-connected}; adding redundant elements
to the basis does not change $\tilde A$; and increasing $N$ changes $A_{(N)}$ only within
the convex neighbourhoods used above. So the class of $\tilde A$ is well defined, and
applying the construction to $X\times[0,1]$ shows it depends only on the homotopy class
of $A$. Reversing the roles, using an orthonormal set of smooth sections of $E$ over $Y$
(which exists in any finite number since $\dim Y>0$) gives the inverse map.
\end{proof}

\subsection{The trace}
For $A\in\Psido{-\infty}(X;E)$ the restriction of the kernel to the diagonal is a
section of $\hom(E)\otimes\Omega$, and we set
\begin{equation}\label{manifolds:eq:trace}
\Tr(A)=\int_X\tr\big(A|_{\Diag}\big)=\int_X\tr k_A(x,x)\,\nu(x).
\end{equation}
For a finite-rank operator $Au=\sum_{j,k}a_{jk}\inner u{v_k}v_j$ with orthonormal $v_j$ as above,
$k_A(x,x)=\sum_{j,k}a_{jk}v_j(x)\otimes v_k(x)^\dagger$, so
$\Tr(A)=\sum_{j,k}a_{jk}\inner{v_j}{v_k}=\sum_ja_{jj}$, the trace of its matrix; since the
matrix trace is invariant under change of basis,
$\tr(b^{-1}ab)=\sum_{i,j,k}b^{-1}_{ij}a_{jk}b_{ki}=\tr a$, this is consistent.

\begin{proposition}\label{manifolds:prop:trace}
The trace \eqref{manifolds:eq:trace} is a continuous linear functional on
$\Psido{-\infty}(X;E)$, and
\begin{enumerate}[label=\textup{(\alph*)}]
\item $\Tr(AB)=\Tr(BA)$ for $A\in\Psido{-\infty}(X;E,F)$, $B\in\Psido{-\infty}(X;F,E)$;
\item each $A\in\Psido{-\infty}(X;E)$ is of trace class on $L^2(X;E)$, with
$\|A\|_1\le C\|k_A\|_{\mathcal C^k}$ for some $k$, and \eqref{manifolds:eq:trace} is its
Hilbert space trace; in particular $\Tr(A)=\lim_N\Tr(A_N)$ for any finite-rank
approximations as in \cref{manifolds:lem:density};
\item $\Tr(AB)=\Tr(BA)$ for $A\in\Psido{-\infty}(X;E)$ and $B\in\Psido m(X;E)$, any $m$.
\end{enumerate}
\end{proposition}
\begin{proof}
Continuity is clear. (a) By Fubini,
\begin{multline*}
\Tr(AB)=\int_X\int_X\tr\big(k_A(x,y)k_B(y,x)\big)\nu(y)\nu(x)
\\=\int_X\int_X\tr\big(k_B(y,x)k_A(x,y)\big)\nu(x)\nu(y)=\Tr(BA),
\end{multline*}
using $\tr(ab)=\tr(ba)$ for $a\in\hom(E_y,F_x)$, $b\in\hom(F_x,E_y)$.

(b) Let $s$ be an even integer with $s>n/2$ and $\Lambda_s$ as in
\cref{manifolds:cor:invertible}. Then $A=\Lambda_s^{-1}(\Lambda_sA)$, where $\Lambda_s^{-1}$ is
Hilbert--Schmidt by \cref{manifolds:lem:HS}, and $\Lambda_sA$ is smoothing with
$\|\Lambda_sA\|_2^2=\int\|\Lambda_sk_A(\cdot,y)\|_{L^2}^2\nu(y)\le C\int\|k_A(\cdot,y)\|^2_{s}\nu(y)\le C'\|k_A\|^2_{\mathcal C^s}$.
So $A$ is of trace class with $\|A\|_1\le\|\Lambda_s^{-1}\|_2\|\Lambda_sA\|_2\le C\|k_A\|_{\mathcal C^s}$
(\cref{hilbert:def:trace-class}). For finite-rank $A$, \eqref{manifolds:eq:finite-rank} gives
Hilbert space trace $\sum_i\inner{f_i}{e_i^\dagger}=\int\sum_i e_i(x)(f_i(x))\nu(x)$, which is
\eqref{manifolds:eq:trace}. Both traces are continuous on $\Psido{-\infty}$ (the first by
the estimate just proved), so they agree by \cref{manifolds:lem:density}.

(c) Write $B=\Lambda_m(\Lambda_m^{-1}B)$ with $\Lambda_m^{-1}B\in\Psido0$ bounded. The operators
$A\Lambda_m$ and $BA$ are smoothing, hence of trace class, and using that
$\Tr(ST)=\Tr(TS)$ for $S$ of trace class and $T$ bounded (\cref{hilbert:prop:trace}),
$\Tr(AB)=\Tr\big((A\Lambda_m)(\Lambda_m^{-1}B)\big)=\Tr(\Lambda_m^{-1}BA\Lambda_m)
=\Tr\big((BA\Lambda_m)\Lambda_m^{-1}\big)=\Tr(BA)$.
\end{proof}

\begin{remark}\label{manifolds:rem:determinant}
Using the trace one can define the Fredholm determinant: there is a unique entire
function $\det(\Id+A)$ on $\Psido{-\infty}(X;E)$ such that
\begin{equation*}
\det\big((\Id+A)(\Id+B)\big)=\det(\Id+A)\det(\Id+B),\qquad
\tfrac{d}{ds}\det(\Id+sA)|_{s=0}=\Tr(A),
\end{equation*}
and $A\in\Gsm(X;E)$ if and only if
$\det(\Id+A)\ne0$. It follows that $\Gsm(X;E)$ is dense in $\Psido{-\infty}(X;E)$. Since
$\Gsm$ is connected, the determinant can be defined on $\Gsm$ by
\begin{equation*}
\det(\Id+A)=\exp\Big(\int_0^1\Tr\big((\Id+\gamma(t))^{-1}\tfrac{d\gamma}{dt}\big)dt\Big)
\end{equation*}
for a smooth curve $\gamma$ in $\Gsm$ from $0$ to $A$; the main points are independence of
the curve and the extension to all of $\Psido{-\infty}$. This is carried out in
\cref{ch:index}; compare also \cref{ch:determinants}.
\end{remark}

\begin{proposition}\label{manifolds:prop:spectrum-smoothing}
Let $A\in\Psido{-\infty}(X;E)$. The set
\begin{equation*}
\Spec(A)\setminus\{0\}=\{z\in\C\setminus\{0\};\ z\Id-A \text{ is not invertible on } L^2(X;E)\}
\end{equation*}
is contained in
$\{|z|\le\|A\|\}$ and is discrete in $\C\setminus\{0\}$; each of its points is an eigenvalue,
and the generalized eigenspace
\begin{equation*}
E(z)=\{u\in L^2(X;E);\ (z\Id-A)^Nu=0 \text{ for some } N\in\N\}
\end{equation*}
is finite-dimensional and contained in $\CI(X;E)$.
\end{proposition}
\begin{proof}
For $|z|>\|A\|$, $z-A=z(\Id-A/z)$ is invertible by the Neumann series. Since $A$ is
compact, the other statements, except smoothness, are the Riesz theory of compact
operators (\cref{hilbert:thm:riesz}). If $(z-A)^Nu=0$ with $z\ne0$, expanding by the binomial
theorem gives $z^Nu=Ap(A)u$ for a polynomial $p$, so $u\in\Ran(A)\subset\CI(X;E)$.
\end{proof}

Alternatively, $\Spec(A)\setminus\{0\}$ is the zero set of the entire function
$z\mapsto\det(\Id-A/z)$ on $\C\setminus\{0\}$, which tends to one as $|z|\to\infty$.

\section{Elliptic complexes and the Hodge theorem}\label{manifolds:sec:hodge}

\subsection{Elliptic complexes}
Let $E_0,E_1,\dots,E_N$ be vector bundles over $X$ and $P_i\in\Diff k(X;E_i,E_{i+1})$
differential operators with $P_{i+1}P_i=0$, so that
\begin{equation}\label{manifolds:eq:complex}
0\longrightarrow\CI(X;E_0)\xrightarrow{P_0}\CI(X;E_1)\xrightarrow{P_1}\cdots
\xrightarrow{P_{N-1}}\CI(X;E_N)\longrightarrow0
\end{equation}
is a complex (we set $P_{-1}=0$, $P_N=0$). Then
$\sigma_k(P_{i+1})\sigma_k(P_i)=\sigma_{2k}(P_{i+1}P_i)=0$, and the complex is called
\emph{elliptic} if the symbol sequence
\begin{equation*}
0\to E_{0,x}\xrightarrow{\sigma_k(P_0)(x,\xi)}E_{1,x}\to\cdots\to E_{N,x}\to0
\end{equation*}
is exact for every $(x,\xi)\in T^*X\setminus0$. Its cohomology is
$H^i=\Nul(P_i)/P_{i-1}\CI(X;E_{i-1})$, where $P_i$ acts on smooth sections.

Choose a positive density and hermitian metrics, and let $P_i^*$ be the formal
adjoints; they form the adjoint complex, $P_{i-1}^*P_i^*=(P_iP_{i-1})^*=0$. Set
\begin{equation*}
\Delta_i=P_i^*P_i+P_{i-1}P_{i-1}^*\in\Diff{2k}(X;E_i).
\end{equation*}

\begin{theorem}[Hodge]\label{manifolds:thm:hodge}
For an elliptic complex \eqref{manifolds:eq:complex} on a compact manifold:
\begin{enumerate}[label=\textup{(\alph*)}]
\item each $\Delta_i$ is elliptic and formally self-adjoint, and its null space
\begin{equation*}
\mathcal H^i=\{u\in\CI(X;E_i);\ P_iu=0,\ P_{i-1}^*u=0\}
\end{equation*}
is finite-dimensional;
\item there is an $L^2$-orthogonal decomposition
\begin{equation}\label{manifolds:eq:hodge-decomposition}
\CI(X;E_i)=\mathcal H^i\oplus P_{i-1}\CI(X;E_{i-1})\oplus P_i^*\CI(X;E_{i+1});
\end{equation}
\item the ranges of the $P_i$ on smooth sections are closed, and we have
$\Nul(P_i)=\mathcal H^i\oplus P_{i-1}\CI(X;E_{i-1})$, so that the natural map
$\mathcal H^i\to H^i$ is an isomorphism; in particular the cohomology of the complex is
finite-dimensional.
\end{enumerate}
\end{theorem}
\begin{proof}
(a) The principal symbol is
$\sigma_{2k}(\Delta_i)=\sigma_k(P_i)^*\sigma_k(P_i)+\sigma_k(P_{i-1})\sigma_k(P_{i-1})^*$, a
self-adjoint endomorphism. If it annihilates $v\in E_{i,x}$, $\xi\ne0$, then
$0=\langle\sigma_{2k}(\Delta_i)v,v\rangle=|\sigma_k(P_i)v|^2+|\sigma_k(P_{i-1})^*v|^2$, so
$v\in\Nul\sigma_k(P_i)=\Ran\sigma_k(P_{i-1})$ by exactness, and also $v\perp\Ran\sigma_k(P_{i-1})$;
hence $v=0$. Being homogeneous of degree $2k$ and invertible on $T^*X\setminus0$, the
symbol is elliptic (\cref{manifolds:lem:elliptic}). For $u\in\CI$,
$\inner{\Delta_iu}u=\|P_iu\|^2+\|P_{i-1}^*u\|^2$, which identifies the null space, and it
is finite-dimensional by \cref{manifolds:thm:fredholm}.

(b) Since $\Delta_i$ is self-adjoint, $\Nul(\Delta_i^*)=\mathcal H^i$, and the generalized
inverse $G_i$ of \cref{manifolds:prop:generalized-inverse} satisfies
$\Delta_iG_i=\Id-\pi_i$, where $\pi_i$ is orthogonal projection onto $\mathcal H^i$. Thus for
$u\in\CI(X;E_i)$
\begin{equation*}
u=\pi_iu+P_{i-1}\big(P_{i-1}^*G_iu\big)+P_i^*\big(P_iG_iu\big).
\end{equation*}
The three spaces are mutually orthogonal: $\inner h{P_{i-1}v}=\inner{P_{i-1}^*h}v=0$ and
$\inner h{P_i^*w}=\inner{P_ih}w=0$ for $h\in\mathcal H^i$, and
$\inner{P_{i-1}v}{P_i^*w}=\inner{P_iP_{i-1}v}w=0$.

(c) If $P_iu=0$, then $w=P_i^*P_iG_iu$ satisfies $P_iw=P_i(u-\pi_iu-P_{i-1}P_{i-1}^*G_iu)=0$,
so $\|w\|^2=\inner{P_iG_iu}{P_iw}=0$. Hence $\Nul(P_i)=\mathcal H^i\oplus P_{i-1}\CI(X;E_{i-1})$,
and $P_{i-1}\CI(X;E_{i-1})=\Nul(P_i)\cap(\mathcal H^i)^\perp$ is closed in $\CI$ (and its closure in
$L^2$ meets $\CI$ in itself). The map $\mathcal H^i\to H^i$ is therefore surjective, and it is
injective since $\mathcal H^i\perp P_{i-1}\CI(X;E_{i-1})$.
\end{proof}

\subsection{The de~Rham complex and the Laplacian on forms}
The exterior derivative $d:\CI(X;\Lambda^p)\to\CI(X;\Lambda^{p+1})$ extends the
differential of functions, satisfies $d^2=0$, and in local coordinates is
$d=\sum_jdx_j\wedge\pa_{x_j}=i\sum_jdx_j\wedge D_{x_j}$. As computed in
\cref{manifolds:sec:definition}, $\sigma_1(d)(x,\xi)=i\,\xi\wedge$. The de~Rham complex
\begin{equation}\label{manifolds:eq:derham}
0\to\CI(X;\Lambda^0)\xrightarrow d\CI(X;\Lambda^1)\xrightarrow d\cdots\xrightarrow d\CI(X;\Lambda^n)\to0
\end{equation}
is elliptic: for $\xi\ne0$, the forms $\alpha\in\Lambda^p_x$ with $\xi\wedge\alpha=0$ are exactly
those of the form $\xi\wedge\beta$, $\beta\in\Lambda^{p-1}_x$, as one sees by choosing a basis
$e_1=\xi,e_2,\dots,e_n$ of $T^*_xX$ and expanding $\alpha$ in the induced basis of
$\Lambda^p_x$. The cohomology of \eqref{manifolds:eq:derham} is the de~Rham cohomology
$H^p_{\mathrm{dR}}(X)$.

Choose a Riemannian metric $g$ on $X$, which induces hermitian metrics on the
$\Lambda^p$ and the Riemannian density. Let $d^*$ be the formal adjoint of $d$ and
\begin{equation*}
\Delta=dd^*+d^*d\in\Diff2(X;\Lambda^p),\qquad\sigma_2(\Delta)(x,\xi)=|\xi|_g^2\Id,
\end{equation*}
since $\sigma_1(d^*)=-i\,\iota_\xi$ (contraction with the dual vector of $\xi$) and
$\xi\wedge\iota_\xi+\iota_\xi\,\xi\wedge=|\xi|^2$. On functions, $\Delta$ is the
Laplace--Beltrami operator, non-negative with null space the locally constant
functions.

\begin{corollary}[Hodge theorem]\label{manifolds:cor:hodge-derham}
On a compact Riemannian manifold, the space $\mathcal H^p$ of harmonic $p$-forms,
$\Delta u=0$ (equivalently $du=0$ and $d^*u=0$), is finite-dimensional, every de~Rham
cohomology class contains a unique harmonic representative, and
\begin{equation*}
\CI(X;\Lambda^p)=\mathcal H^p\oplus d\,\CI(X;\Lambda^{p-1})\oplus d^*\CI(X;\Lambda^{p+1}).
\end{equation*}
\end{corollary}

Thus the Betti numbers $b_p=\dim H^p_{\mathrm{dR}}(X)$ are finite. On the full form
bundle $\Lambda^*=\bigoplus_p\Lambda^p$, the operator $d+d^*$ satisfies
$(d+d^*)^2=\Delta$, so is elliptic, and it maps even forms to odd forms.

\begin{corollary}\label{manifolds:cor:euler}
The index of $d+d^*:\CI(X;\Lambda^{\mathrm{ev}})\to\CI(X;\Lambda^{\mathrm{odd}})$ is the Euler
characteristic $\chi(X)=\sum_p(-1)^pb_p$.
\end{corollary}
\begin{proof}
If $(d+d^*)u=0$ then $\Delta u=0$; conversely if $\Delta u=0$ then each component of $u$
satisfies $du_p=0$, $d^*u_p=0$. So the null space of $d+d^*$ on even forms is
$\bigoplus_p\mathcal H^{2p}$; its adjoint is $d+d^*$ on odd forms, with null space
$\bigoplus_p\mathcal H^{2p+1}$. Now apply \cref{manifolds:thm:fredholm}(d) and
\cref{manifolds:cor:hodge-derham}.
\end{proof}

In particular, the index of this operator does not depend on the Riemannian metric,
as also follows from the homotopy invariance of the index.

\section{Pseudodifferential projections}\label{manifolds:sec:projections}
In this section $E$ is a hermitian bundle over the compact manifold $X$, with $L^2$
defined by a positive density.

\begin{proposition}\label{manifolds:prop:projection}
If $P\in\Psido0(X;E)$ satisfies $P^2-P\in\Psido{-\infty}(X;E)$ then there exists
$\Pi\in\Psido0(X;E)$ with $\Pi^2=\Pi$ and $\Pi-P\in\Psido{-\infty}(X;E)$.
\end{proposition}
\begin{proof}
Let $R=P^2-P$, which commutes with $P$. For $\lambda\notin\{0,1\}$ a direct computation gives
\begin{equation}\label{manifolds:eq:projection-resolvent}
(\lambda-P)(P+\lambda-1)=\lambda(\lambda-1)-R=\lambda(\lambda-1)\big(\Id-S_\lambda\big),\qquad
S_\lambda=\frac R{\lambda(\lambda-1)},
\end{equation}
and all these operators commute. Hence if $\lambda(\lambda-1)\notin\Spec(R)$, then
$\lambda-P$ is invertible, with
\begin{equation*}
(\lambda-P)^{-1}=\frac{P+\lambda-1}{\lambda(\lambda-1)}(\Id-S_\lambda)^{-1}
=\frac{P+\lambda-1}{\lambda(\lambda-1)}+\frac{P+\lambda-1}{\lambda(\lambda-1)}\,S_\lambda(\Id-S_\lambda)^{-1}.
\end{equation*}
Since $R$ is compact, $\Spec(R)\setminus\{0\}$ is countable (\cref{hilbert:thm:riesz}), so the
set of $r\in(\frac14,\frac34)$ such that the circle $\gamma_r=\{|\lambda-1|=r\}$ meets
$\{\lambda;\lambda(\lambda-1)\in\Spec(R)\}$ is countable, and we may fix $r$ with $\gamma=\gamma_r$
disjoint from $\Spec(P)$. Let
\begin{equation*}
\Pi=\frac1{2\pi i}\oint_\gamma(\lambda-P)^{-1}d\lambda .
\end{equation*}
The first term in the expansion of the resolvent contributes, by the residue
theorem (the only pole inside $\gamma$ is at $\lambda=1$), exactly $P$. In the second term,
$S_\lambda(\Id-S_\lambda)^{-1}=S_\lambda+S_\lambda(\Id-S_\lambda)^{-1}S_\lambda$ is smoothing, with kernel
depending continuously on $\lambda\in\gamma$ by \cref{manifolds:lem:sandwich}, and remains so
after composition with $P+\lambda-1$; so its integral is smoothing. Hence
$\Pi-P\in\Psido{-\infty}$.

That $\Pi^2=\Pi$ is the standard property of Riesz projections: with $\gamma'=\gamma_{r'}$,
$r'<r$ close to $r$ and also disjoint from $\Spec(P)$, the resolvent identity
$(\lambda-P)^{-1}(\mu-P)^{-1}=\big((\lambda-P)^{-1}-(\mu-P)^{-1}\big)/(\mu-\lambda)$ gives
\begin{equation*}
\Pi^2=\frac1{(2\pi i)^2}\oint_{\gamma'}\oint_\gamma\frac{(\lambda-P)^{-1}-(\mu-P)^{-1}}{\mu-\lambda}\,d\mu\,d\lambda
=\frac1{2\pi i}\oint_{\gamma'}(\lambda-P)^{-1}d\lambda=\Pi,
\end{equation*}
since $\oint_\gamma(\mu-\lambda)^{-1}d\mu=2\pi i$ for $\lambda$ on $\gamma'$ (inside $\gamma$) and
$\oint_{\gamma'}(\mu-\lambda)^{-1}d\lambda=0$ for $\mu$ on $\gamma$; the integral over $\gamma'$ equals that
over $\gamma$ by Cauchy's theorem, the resolvent being holomorphic between the circles.
\end{proof}

\begin{proposition}\label{manifolds:prop:range-projection}
Let $P\in\Psido0(X;E)$ satisfy $P^2-P\in\Psido{-\infty}(X;E)$, and let $F\subset\Hs s(X;E)$ be a
closed subspace for which there are $A,B\in\Psido{-\infty}(X;E)$ with
\begin{equation*}
(\Id-P)u=Au\ \ \forall u\in F\quad\text{and}\quad(P+B)\Hs s(X;E)\subset F.
\end{equation*}
Then there exists $B'\in\Psido{-\infty}(X;E)$ such that $(P+B')^2=P+B'$ and
$F=(P+B')\Hs s(X;E)$.
\end{proposition}
\begin{proof}
Assume first that $s=0$. All modifications below change $P$ by smoothing operators,
and the hypotheses are preserved, with new $A$ and $B$, under such changes: if
$P'=P+C$ with $C$ smoothing, then $\Id-P'=A-C$ on $F$ and $P'+(B-C)=P+B$.

\emph{Step 1.} By \cref{manifolds:prop:projection} we may assume that $P$ is a projection,
$P^2=P$. Then $\Ran(\Id-P)=\Nul(P)$.

\emph{Step 2.} Let $E_0=F\cap\Nul(P)$. On $E_0$, $u=(\Id-P)u=Au$, so $E_0\subset\Nul(\Id-A)$,
which is finite-dimensional and consists of smooth sections
(\cref{manifolds:thm:fredholm}). Let $e_1,\dots,e_k$ be an orthonormal basis of $E_0$ and
$f_j=(\Id-P^*)e_j\in\CI(X;E)$; then $\inner{e_l}{f_j}=\inner{(\Id-P)e_l}{e_j}=\delta_{lj}$, so
$Qu=\sum_j\inner u{f_j}e_j$ is a smoothing projection onto $E_0$. Moreover $QP=0$, since
$P^*f_j=0$, and $PQ=0$, since $e_j\in\Nul(P)$; hence $P_1=P+Q$ is a projection. If $u\in F$
and $P_1u=0$, then $Pu=-Qu\in E_0\subset\Nul(P)$, so $Pu=P^2u=0$, then $Qu=0$, $u\in E_0$ and
$u=Qu=0$. Thus, replacing $P$ by $P_1$, we may assume $F\cap\Nul(P)=\{0\}$.

\emph{Step 3.} Let $G=F+\Nul(P)$. Since $\Id+B=(P+B)+(\Id-P)$, the range of $\Id+B$ lies
in $G$. By \cref{manifolds:thm:fredholm}, $\Ran(\Id+B)$ is closed, with $L^2$-orthogonal
complement $N'=\Nul(\Id+B^*)\subset\CI(X;E)$ of finite dimension. Hence
$G=\Ran(\Id+B)\oplus(G\cap N')$ is closed, and if $C'$ is the orthocomplement of $G\cap N'$
in $N'$ then $L^2(X;E)=G\oplus C'$ with $C'\subset\CI$ finite-dimensional. Let
$c_1,\dots,c_l$ be a basis of $C'$. Since $C'\cap\Nul(P)=\{0\}$, the $c'_j=Pc_j$ are linearly
independent; let $h_j$ be the dual basis of their span with respect to the $L^2$ inner product
and $g_j=P^*h_j$. Then $Q'u=\sum_j\inner u{g_j}c'_j$ is a smoothing operator with $Q'P=Q'$,
$PQ'=Q'$, $Q'c'_j=c'_j$, so $Q'^2=Q'$ and $P_2=P-Q'$ is a projection. One checks that
$\Nul(P_2)=\Nul(P)\oplus C'$: indeed $P_2c_j=c'_j-Q'Pc_j=0$, $\Nul(P)\subset\Nul(P_2)$, and if
$P_2u=0$ then $Pu=Q'u=\sum a_jPc_j$, so $u-\sum a_jc_j\in\Nul(P)$. Consequently
$L^2=F\oplus\Nul(P_2)$: the sum is $G\oplus C'$, and if $f=v+c$ with $f\in F$, $v\in\Nul(P)$,
$c\in C'$, then $c\in G\cap C'=\{0\}$ and $f\in F\cap\Nul(P)=\{0\}$. Replacing $P$ by $P_2$, we may
therefore assume
\begin{equation}\label{manifolds:eq:direct-sum}
L^2(X;E)=F\oplus\Nul(P),
\end{equation}
a direct, not necessarily orthogonal, sum of closed subspaces.

\emph{Step 4.} Let $P'$ be the projection with range $F$ and null space $\Nul(P)$, which is
bounded by \eqref{manifolds:eq:direct-sum} and the closed graph theorem. Since
$\Nul(P')=\Nul(P)=\Ran(\Id-P)$, we have $P'(\Id-P)=0$ and $P(\Id-P')=0$, so $P'=P'P$ and
$P=PP'$. Hence
\begin{equation*}
P'-P=(\Id-P)P'=AP',
\end{equation*}
since $P'u\in F$ and $\Id-P=A$ on $F$. Also $P'(P+B)=P+B$ since $(P+B)L^2\subset F$, so
$P'-P=P'P-P=B-P'B=(\Id-P')B$. Combining,
\begin{equation*}
P'-P=AP'=A\big(P+(\Id-P')B\big)=AP+A(\Id-P')B,
\end{equation*}
which is smoothing by \cref{manifolds:cor:smoothing-ideal} and \cref{manifolds:lem:sandwich}.
Thus $B'=P'-P$ (relative to the original $P$, plus the smoothing changes made in
Steps 1--3) has the required properties.

For general $s$, let $\Lambda_s$ be as in \cref{manifolds:cor:invertible}. Then
$\tilde P=\Lambda_sP\Lambda_s^{-1}\in\Psido0$ satisfies $\tilde P^2-\tilde P\in\Psido{-\infty}$,
$\tilde F=\Lambda_sF\subset L^2$ is closed, and the hypotheses hold with $\tilde A=\Lambda_sA\Lambda_s^{-1}$,
$\tilde B=\Lambda_sB\Lambda_s^{-1}$. The case $s=0$ gives $\tilde B'$, and
$B'=\Lambda_s^{-1}\tilde B'\Lambda_s$ has the required properties.
\end{proof}

\begin{remark}\label{manifolds:rem:functional-calculus}
The functional calculus of elliptic operators, through the heat kernel
$e^{-tA}$, the resolvent $(A-\lambda)^{-1}$ with parameter, and the complex powers $A^z$
of a positive elliptic operator, is not developed here; for these see
\cite{Seeley} and \cite{Shubin}.
\end{remark}

\section{Exercises}\label{manifolds:sec:exercises}

\begin{exercise}\label{manifolds:ex:compatible}
Show that compatibility, as defined in \cref{manifolds:def:structure}, is an
equivalence relation on $\CI$ structures. Conclude that there is a unique maximal
$\CI$ structure containing any given $\CI$ structure.
\end{exercise}

\begin{exercise}\label{manifolds:ex:maximal}
Let $\mathcal F$ be a $\CI$ structure on $X$ and let $O_a$, $a\in\mathcal A$, be a covering of
$X$ by coordinate neighbourhoods as in \eqref{manifolds:eq:coord}. Show that the maximal
$\CI$ structure containing $\mathcal F$ consists of all functions on $X$ which, on each
$O_a$, are of the form $g\circ\kappa_a$ with $g$ smooth. Conclude that the maximal $\CI$
structure is an algebra.
\end{exercise}

\begin{exercise}[Partitions of unity]\label{manifolds:ex:partition}
Show that any $\CI$ manifold admits partitions of unity. That is, if $O_a$, $a\in\mathcal A$,
is an open cover of $X$, then there exist $\rho_{a,i}\in\CI(X)$, $a\in\mathcal A$, $i\in\N$, with
$0\le\rho_{a,i}\le1$, each $\rho_{a,i}$ vanishing outside a compact subset $K_{a,i}\subset O_a$,
such that each point has a neighbourhood meeting only finitely many of the $K_{a,i}$,
and $\sum_{a,i}\rho_{a,i}=1$.
\end{exercise}

\begin{exercise}\label{manifolds:ex:restriction}
Fill in the details in the proof of \cref{manifolds:prop:restriction}: with $u_R$
obtained by replacing $a$ by $a(w',t,\zeta)\chi(\zeta/R)$, show that $u_R\to u$ in $\CmI(M)$ and
that $u_R|_Z$ converges in $\CmI(Z)$ as $R\to\infty$. Deduce that $u|_Z$ does not depend on
the coordinates used, and that $(\phi u)|_Z=\phi|_Z\,u|_Z$ for $\phi\in\CI(M)$.
\end{exercise}

\begin{exercise}\label{manifolds:ex:smoothing-char}
Show that a continuous linear map $A:\CI(X;E)\to\CI(X;F)$ is smoothing if and only if
it extends to a continuous linear map $\CmI(X;E)\to\CI(X;F)$. Show that smoothing operators are \emph{not}
characterized as the bounded operators $L^2(X)\to\CI(X)$, by considering
$Au=\inner ug\cdot1$ with $g\in L^2(X)\setminus\CI(X)$.
\end{exercise}

\begin{exercise}\label{manifolds:ex:oscillatory}
Let $A\in\Psido m_{\mathrm{ph}}(X;E,F)$ and let $f\in\CI(X)$ be real-valued with
$df(\bar x)\ne0$. Show that for $u\in\CI(X;E)$, as $t\to+\infty$,
$e^{-itf}A(e^{itf}u)(\bar x)=t^m\sigma_m(A)(\bar x,df(\bar x))u(\bar x)+O(t^{m-1})$, where
$\sigma_m(A)$ is the homogeneous principal symbol. (Use the left-reduced form in local
coordinates and Taylor's formula for $f$.) Conclude that this gives an invariant
definition of the principal symbol of classical operators extending
\eqref{manifolds:eq:oscillatory-test}.
\end{exercise}

\begin{exercise}\label{manifolds:ex:index-zero}
Show that an elliptic $A\in\Psido m(X;E)$ which is formally self-adjoint has index zero,
and that the same is true of any elliptic operator in $\Psido m(X;E)$ whose principal
symbol can be connected to that of a self-adjoint elliptic operator through elliptic
symbols.
\end{exercise}

\begin{exercise}\label{manifolds:ex:poincare}
Let $X$ be compact, oriented and Riemannian, of dimension $n$, and let
$*:\Lambda^p\to\Lambda^{n-p}$ be the Hodge star. Show that $d^*=\pm*d*$ on $p$-forms and that
$*$ maps harmonic $p$-forms isomorphically onto harmonic $(n-p)$-forms. Deduce Poincar\'e
duality $b_p=b_{n-p}$, and that $\chi(X)=0$ if $n$ is odd.
\end{exercise}

\begin{exercise}\label{manifolds:ex:G-dense}
Assuming the properties of the Fredholm determinant in
\cref{manifolds:rem:determinant}, show that $\Gsm(X;E)$ is dense in
$\Psido{-\infty}(X;E)$. Prove this directly using \cref{manifolds:lem:density} and the
density of $\GL(M,\C)$ in $M(M,\C)$.
\end{exercise}

\section*{Notes}
This chapter combines Chapter~5 of the 2003 notes with Lectures~5--8 of the 2005
lectures. The invariant definition of pseudodifferential operators through
conormal kernels comes from the latter, and the local definition from the
former; \cref{manifolds:thm:local} shows that they agree. The source text gave
only a sketch of the Fredholm property, carried out on $\CI$ with $L^2$
estimates; here it is proved on Sobolev spaces, which are defined in
\cref{manifolds:sec:sobolev}, and the range is characterized by duality. The Hodge
theorem, only announced in the 2003 notes, is proved for general elliptic
complexes as in the 2005 lectures. The existence of the projection $\Pi$ in
\cref{manifolds:prop:projection} was stated without proof; we use a contour integral.
The functional calculus (heat kernel, resolvent, complex powers) announced in the
original is not included; see \cite{Seeley,Shubin}. For a full treatment of
pseudodifferential operators on manifolds see \cite[Chapter~18]{Hormander-III}
and \cite{Shubin}.
